% =====================================================================
% Appendix for main.tex. Sections carry the labels referenced from the
% main text: app:proofs, app:support, app:lean, app:protocol, app:eeg,
% app:hyper, app:repro.
% =====================================================================

% Compact restatements of main-text results.
\newenvironment{restated}[1]{%
  \par\medskip\noindent\textbf{#1 (restated).}\hspace{0.4em}\ignorespaces
}{\par\smallskip}

\section{Proofs of Section~\ref{sec:theory}}
\label{app:proofs}

This appendix proves the results of Section~\ref{sec:theory} in full, restating each numbered claim or setting it up in the notation of its proof.
Throughout, $(x_t, y_t)_{t \ge 1}$ is a stream with $x_t \in \R^D$ obeying
$y_t = x_t^\top w^*_t + \varepsilon_t$, where $\mathcal F_t$ is the $\sigma$-algebra generated by the past, $\E[\varepsilon_t \mid \mathcal F_t] = 0$, and $\E[\varepsilon_t^2 \mid \mathcal F_t] \le \sigma^2$.
The learner maintains sufficient statistics
\[
  A_t = \alpha A_{t-1} + x_t x_t^\top, \qquad
  B_t = \alpha B_{t-1} + x_t y_t, \qquad
  n_t = \alpha n_{t-1} + 1,
\]
with forgetting $\alpha \in (0,1]$ and initialization from a warm-up set of $T_w$ samples.
Write $\bar A_t = A_t/n_t$ with eigendecomposition $\bar A_t = U\,\mathrm{diag}(\mu_1,\dots,\mu_D)\,U^\top$ and $\bar b_t = B_t/n_t$; an index $i$ always refers to the coordinate in this eigenbasis.
\emph{Two design models.} Lemma~\ref{lem:zerolag} and Theorems~\ref{thm:cadence}--\ref{thm:window} are proved in the \emph{mean-field} model $\bar A_t \equiv \Sigma$, i.e.\ the random accumulator is replaced by its population expectation; this is the model in which the lag, the variance factor $\Psi$, and the window target are exact. Theorem~\ref{thm:rmt} is the finite-$D$ \emph{random-design} correction: it retains the random $\bar A_t$ and takes the proportional limit $\zeta = D\omega$ fixed. The two agree as $\zeta \to 0$; the grid's pre-drift windows have $\zeta \le 0.13$, while drift-phase windows reach $\zeta \le 6.4$, where the exact random-design variance of Theorem~\ref{thm:rmt} replaces the mean-field leading-order term so the cycle-averaged optimum is available across the whole range (Appendix~\ref{app:proofs}, ``Exact finite-window variance term'').
The anchored solve with anchor $G_t$ and strength $\rho > 0$ is $W_t = (A_t + n_t\rho I)^{-1}(B_t + n_t\rho G_t)$; in the eigenbasis, coordinate $i$ is the scalar problem
\begin{equation}
\label{eq:app-anchored}
  W_{t,i} \;=\; \frac{\bar b_{t,i} + \rho\,G_{t,i}}{\mu_{t,i} + \rho}
  \;=\; (1 - s_i)\,\widehat w_{t,i} + s_i\,G_{t,i},
  \qquad
  s_i \;=\; \frac{\rho}{\mu_{t,i} + \rho},
\end{equation}
where $\widehat w_{t,i} = \bar b_{t,i}/\mu_{t,i}$ is the per-direction data solution.
Under the warm-up model the data solution is unbiased per direction with variance $v_i = \sigma^2/(n_t \mu_{t,i})$.
Two parametrizations of the anchored family are used: the \emph{scalar-$\rho$} family (one $\rho$; the pull $s_i = \rho/(\mu_i + \rho)$ is then spectrum-adaptive) and the \emph{uniform-$s$} family (one pull $s$ for all directions, i.e.\ $\rho_i = \mu_i s/(1-s)$; the solution path is the segment between $\widehat w$ and $G$).
Both satisfy the pull identity, which underlies every result below.

\subsection{Proof of Equation~\eqref{eq:pull} (pull identity)}
\label{app:proof-pull}



\begin{proof}[Proof of Equation~\eqref{eq:pull}]
Since $1 - s = \mu/(\mu + \rho)$,
\[
  (1 - s)\,\widehat w + s\,\gamma g
  \;=\; \frac{\mu}{\mu + \rho}\cdot\frac{b}{\mu}
        + \frac{\rho}{\mu + \rho}\cdot \gamma g
  \;=\; \frac{b + \rho \gamma g}{\mu + \rho}
  \;=\; W,
\]
and subtracting $w = (1-s)w + s w$ gives the identity.
No approximation is involved.
Applying it coordinatewise in the eigenbasis to \eqref{eq:app-anchored} yields the form stated in Section~\ref{sec:theory}.
Both the $\gamma = 1$ and general-$\gamma$ forms are formalized (\texttt{pull\_identity}, \texttt{pull\_identity\_general}), together with the eigenbasis decoupling of the anchored solve (\texttt{eigenbasis\_solve\_coord}, \texttt{eigenbasis\_pull\_identity}) in \texttt{Core.lean}.
\end{proof}

\subsection{Proof of Lemma~\ref{lem:zerolag} (zero-lag refresh equilibrium)}
\label{app:proof-tracking}



\begin{proof}[Proof of Lemma~\ref{lem:zerolag}]
(i) At equilibrium,
\[
  \E[B_t]
  \;=\; \sum_{k \ge 0} \alpha^k\, \mu\,\bigl(w^*_t - (k+1)\nu\bigr)
  \;=\; \mu\,\Bigl(w^*_t - \frac{\nu}{1-\alpha}\Bigr)\cdot\frac{1}{1-\alpha},
\]
so $\widehat w - w^* = -\nu/(1-\alpha)$ exactly at the mean-field expectation ($\widehat w = \bar b/\mu$ is the unanchored data solve and carries no $\rho$ dependence; the centered estimation noise contributes only the variance terms analyzed in Proposition~\ref{prop:conflict}).
(ii) By the pull identity with $\widehat w - w^* = -\nu/(1-\alpha)$ and $G - w^* = W - w^*$ at the continuous-refresh fixed point:
\[
  W - w^* \;=\; (1 - s)\Bigl(-\frac{\nu}{1-\alpha}\Bigr) + s\,(W - w^*),
\]
hence $(1 - s)(W - w^*) = -(1-s)\nu/(1-\alpha)$, so $W - w^* = -\nu/(1-\alpha)$ for any $s < 1$.
This algebraic step is formalized (\texttt{lag\_invariance} in \texttt{Core.lean}).
(iii) Theorem~\ref{thm:cadence}'s moving-frame recursion gives the exact cycle-averaged lag $\Delta_k(s) = s\nu k\bigl(a/(1-a)+\tfrac12\bigr)$ with $a=s(2-s)$; since $1-a=(1-s)^2$, this is $O\bigl(k s/(1-s)^2\bigr)\nu$.
The prediction-risk cost of the extra lag is $\mu_i\,\Delta_k(s)^2$, which vanishes as $\mu_i \to 0$ at fixed $s < 1$.
\end{proof}

At the continuous-refresh equilibrium the pull identity forces $W = \widehat w$ for $s < 1$: the anchor is inert for mean tracking, so it contributes neither lag nor variance reduction in that limit.

\subsection{Finite-cadence refresh equilibrium (Theorem~\ref{thm:cadence})}
\label{app:cadence}

Work per eigen-direction with pull $s \in (0,1)$ and refresh cadence $k$; let $t_m = mk$ be the refresh steps, $w^*(t)$ the moving target (velocity $\nu$ per step), and $L = \nu/(1-\alpha)$ the pure-leaky lag of the data path (Lemma~\ref{lem:zerolag}(i)). The refresh reads the post-update solve, so the lag entering the recursion is the post-update value $\nu(1-\alpha)/\alpha$, one step shorter than the predictor lag; the cycle-averaged lag relative to the pure-leaky path, $\Delta_k(s)$ below, is unchanged by this one-step offset.
Between refreshes the anchor $G_m$ is frozen, and the anchored solution at $t \in [t_m, t_{m+1})$ is, in the moving frame,
\[
  W_t - w^*(t) \;=\; (1-s)\,(-L) \;+\; s\,\bigl(g_m - \nu\,(t - t_m)\bigr),
  \qquad g_m := G_m - w^*(t_m),
\]
where $\nu (t - t_m)$ is the target's advance since the last refresh, invisible to the frozen anchor.
At the refresh, $G_{m+1} = (1-s)\,W_{t_{m+1}} + s\,G_m$ measured against the new target position, which yields the affine recursion
\[
  g_{m+1} \;=\; a\,g_m - (1-s)^2\,L - a\,\nu k,
  \qquad a := (1-s)s + s = s(2-s) = 1 - (1-s)^2,
\]
with $a \in (0,1)$ for every $s \in (0,1)$, so the recursion contracts at rate $a$ per cycle to
\[
  g^* \;=\; -L - \nu k\,\frac{a}{1-a},
  \qquad
  g_m \;=\; g^* + a^m\,(g_0 - g^*).
\]
The rate identity, the fixed-point equation, the closed form, and the contraction are formalized (\texttt{refresh\_rate\_identity}, \texttt{refresh\_fixed\_point},\newline\texttt{refresh\_closed\_form}, \texttt{refresh\_contraction\_rate}, together with the moving-frame reduction and the midpoint cycle assembly, in \texttt{Core.lean}); we also verified the fixed point and the cycle-averaged lag against a one-dimensional simulation at two operating points.
The geometric steady-lag form parallels the iterate-averaging analyses under drift \citep{polyak1992,cutler2023} and the per-direction momentum lag analysis of \citet{hong2026}.
The cycle-averaged solution lags the pure-leaky path by
\[
  \Delta_k(s) \;=\; s\,\nu k \Bigl(\frac{a}{1-a} + \frac{1}{2}\Bigr)
  \;=\; \frac{s\,\nu k\,(1 + 2s - s^2)}{2\,(1-s)^2},
\]
linear in the cadence and vanishing as $s \to 0$; the within-cycle residual adds a variance term of order $s^2 \nu^2 k^2/12$.
The first summand is the anchor's accumulated staleness---each cycle's drift persists into the anchor with weight $a^j$, summing to $a/(1-a)$---and the second is the within-cycle advance, which makes the within-cycle correction exact and adds the geometric term.
The continuous-refresh limit of Lemma~\ref{lem:zerolag}(ii) is the instantaneous version of the same equilibrium: at the refresh instant the anchor adds no lag beyond $L$, and the correction $\nu k\,a/(1-a)$ is precisely what the target's advance between refreshes contributes.

The data-path noise does not decorrelate between refreshes: the leaky accumulator gives $\mathbb{E}[\epsilon_m \epsilon_{m-j}] = r_k^{\,j} v$ with $r_k = \alpha^k$ ($\alpha = 1-\omega$ the forgetting factor), which at the nominal constants ranges over $r_{10} \in [0.60, 0.995]$. The correlation is material at these operating points.
Writing $\tilde g_m$ for the zero-mean anchor fluctuation, the refresh rule gives the stochastic recurrence $\tilde g_m = (1-s)^2 \epsilon_m + a\,\tilde g_{m-1}$, whose steady-state cross-covariance against the fresh noise is the geometric sum
\[
  C(s,r_k) \;:=\; \mathbb{E}[\epsilon_m \tilde g_{m-1}]
  \;=\; (1-s)^2 \sum_{j \ge 0} a^j r_k^{\,j+1} v
  \;=\; \frac{(1-s)^2\,r_k}{1 - ar_k}\,v .
\]
Solving the steady-state variance system $q = (1-s)^4 v + a^2 q + 2a(1-s)^2 C$ then gives the exact pair
\[
  q(s,r_k) \;=\; \frac{(1-s)^2}{1+a}\cdot\frac{1+ar_k}{1-ar_k}\,v,
\]
\[
  \mathrm{Var}(W) \;=\; (1-s)^2 v + s^2 q + 2s(1-s)\,C
  \;=\; v\,\Bigl(1 - \frac{2s\,(1+s-s^2)\,(1-r_k)}{(1+a)\,(1-ar_k)}\Bigr),
\]
strictly below $v$ for every $s \in (0,1)$ and $r_k \in [0,1)$ (\texttt{refresh\_anchor\_variance},\newline\texttt{refresh\_variance\_form}, \texttt{refresh\_variance\_lt} in \texttt{Core.lean}); we verified the identities symbolically and the refresh-instant variance against one-dimensional simulations at six operating points (deviation $\le 1\%$); this is the variance at the refresh step, and the cycle-averaged variance is within a few percent of it at the operating pulls.
The lower endpoint is exact as well: subtracting $(1-s)^2$ from the bracket factorises as
\[
  \Psi(s,r_k)-(1-s)^2
  = \frac{s\,(1-s)^2\bigl(r_k\,(s^3-4s^2+2s+2)+s\bigr)}{(1+a)\,(1-ar_k)},
\]
and $s^3-4s^2+2s+2 = 1+(s-1)(s^2-3s-1) \ge 1$ on $[0,1]$, so the numerator is nonnegative and $\Psi(s,r_k)\ge(1-s)^2$ (\texttt{refresh\_variance\_lb} in \texttt{Core.lean}).
Hence the frozen-anchor plug-in $c'$ lower-bounds the exact variance and the window estimator errs toward larger, safer windows.
Both endpoints are exact: $r_k \to 0$ (refresh cycles long against the noise memory) recovers the uncorrelated reduction $v[1 - 2s(1+s-s^2)/(1+a)]$, and $r_k \to 1$ (continuous refresh) gives $\mathrm{Var}(W) \to v$---the inert-anchor limit of Lemma~\ref{lem:zerolag}(ii), now an exact endpoint of a single formula rather than a separate statement.
The cadence therefore trades two real costs: refreshing fast lets the anchor co-move with the data path and erases its variance benefit, refreshing slowly inflates the lag term $\mu_i\,\Delta_k(s)^2$, and the default $k = 10$ falls within the tradeoff.
The two costs are paid in different directions: the excess lag blows up as $s \to 1$, but the directions with $s \approx 1$ are precisely the weakly excited ones ($\mu_i \ll \rho$), where the lag's risk contribution is negligible---the filter-bank property of the $\beta_i = 1 - s_i$ refresh, and the reason Proposition~\ref{prop:conflict} drives $\rho^*$ down, and the lag with it, exactly when measured conflict grows.

\medskip
\noindent\emph{Remark (drift-induced cross term).}
Under nonzero drift velocity the risk decomposition $R(s) = (1-s)^2 v + s^2 c$ of Proposition~\ref{prop:conflict} acquires the cross term $2(1-s)s\,\E[(\widehat w - w^*)(G - w^*)]$ between the lagged data solve and the anchor error.
Cauchy--Schwarz bounds it through the two quantities the theory already controls: $|\E[(\widehat w - w^*)(G - w^*)]| \le \sqrt{\E[(\widehat w - w^*)^2]\,\E[(G - w^*)^2]}$, where the first factor is the lag-induced gap (bounded per direction by Theorem~\ref{thm:cadence}'s equilibrium: the moving-average gap $\bar\Delta_k \le (1-s)^2\nu k/(1-a)$ at cadence $k$) and the second is the anchor conflict $c$.
The cross term therefore inflates the effective conflict to $c_{\mathrm{eff}} = c + 2(1-s)s\sqrt{\bar\Delta_k\,c}$ at most, and Proposition~\ref{prop:minimax} applies unchanged: overestimating $c_{\mathrm{eff}}$ remains minimax-safe, so the calibration rule $\rho^* = \bar\mu \widehat V / \widehat C$ with the conservative conflict estimate retains its safety property under drift.

\subsection{Proof of Theorem~\ref{thm:window} (optimal forgetting window)}
\label{app:proof-window}



\begin{proof}[Proof of Theorem~\ref{thm:window}]
\emph{Step 0 (finite cadence to the surrogate).}
The finite-cadence risk \eqref{eq:app-window-risk} combines Theorem~\ref{thm:cadence}'s two exact steady-state objects: the per-direction variance is $v_i\Psi(s_i,r_k)$ with $v_i=\sigma^2/(n_{\mathrm{eff}}\mu_i)$, and the per-direction lag is $\nu_i/\omega+\Delta_k(s_i)$; the within-cycle fluctuation contributes the $O(k^2\nu^2s^2)$ term derived in Appendix~\ref{app:cadence}.
Because $(1-s_i)^2 \le \Psi(s_i,r_k) \le 1$ for all $r_k \in [0,1)$ (\texttt{refresh\_variance\_lb}), the frozen-anchor surrogate $c'\,\omega/(2-\omega)+\nu^2_\Sigma/\omega^2$ lower-bounds the exact cycle-averaged variance term at every cadence, and $D(1-\bar s)^2 \le d_{\mathrm{eff}}(\rho)$ by Jensen; the runtime plug-in therefore underestimates the variance coefficient, returning a window no smaller than the exact-optimal one---the over-forgetting, safety-asymmetric direction of Step~4.
The finite-cadence lag residual is bounded at the runtime range: with $b(s)=s(2-s)/(1-s)^2$ and $k\omega \le 1/2$, $\Delta_k(s_i)=s_i\nu_i k\bigl(b(s_i)+\tfrac12\bigr)\le \tfrac12 s_i\bigl(b(s_i)+\tfrac12\bigr)\nu_i/\omega$, a vanishing fraction of the leading lag ($\le 0.4\%$ for $s_i\le 0.1$ at $k\omega=0.5$) whose prediction cost $\mu_i\Delta_k(s_i)^2$ is concentrated in weakly excited directions.
Numerically, over $k\omega \in [0.005,0.5]$ the surrogate's cubic minimizer is within $2.1\%$ of the exact cycle-averaged optimum for $s_i \le 0.2$ ($5.2\%$ for $s_i \le 0.3$), with the largest gaps at the window floor (\texttt{diagnostics.py --diag subopt}; \texttt{results/window\_subopt.csv}).
In the quasi-static limit $1 \ll k\omega \ll 1/s$, $r_k=(1-\omega)^k \le e^{-\omega k} \to 0$, so $\Psi(s_i,r_k)=(1-s_i)^2(1+O(s_i^2+r_k))$ and $\Delta_k(s_i)=o(\nu_i/\omega)$; both corrections vanish and the reduced form \eqref{eq:app-window-risk} is exact, with remainder $o(1)$ uniformly over the regime.

\emph{Step 1 (risk decomposition).}
For exponential weights $\alpha^k$ the effective sample count of the accumulator is
\[
  n_{\mathrm{eff}}
  \;=\; \frac{\bigl(\sum_{k \ge 0} \alpha^k\bigr)^2}{\sum_{k \ge 0} \alpha^{2k}}
  \;=\; \frac{(1/(1-\alpha))^2}{1/(1-\alpha^2)}
  \;=\; \frac{1 + \alpha}{1 - \alpha}
  \;=\; \frac{2-\omega}{\omega},
\]
so the estimation-variance contribution to the solve risk is $(1-\bar s)^2 \sigma^2 D/n_{\mathrm{eff}} = c'\,\omega/(2-\omega)$.
The tracking contribution is the lag of Lemma~\ref{lem:zerolag}(i) below: per direction the solve lags the target by $\nu_i/\omega$, so the prediction-units lag risk is $\sum_i \mu_i (\nu_i/\omega)^2 = \nu^2_\Sigma/\omega^2$.
Summing the two contributions gives
\begin{equation}
\label{eq:app-window-risk}
  R_k(\omega,s) \;=\; c'\,\frac{\omega}{2-\omega} + \frac{\nu^2_\Sigma}{\omega^2} + o(1),
  \qquad c' = \sigma^2\,d_{\mathrm{eff}}(\rho),
\end{equation}
to leading order in $\omega$, $R(\omega) = c'\omega/2 + \nu^2_\Sigma/\omega^2 + O(\omega^2)$.
The exact random-design refinement of the variance term is $\sigma^2\Upsilon(\omega;\rho)/(1-\Upsilon(\omega;\rho))$ (Theorem~\ref{thm:rmt}(i)), equal to it as $\zeta \to 0$ and accurate to within $5\%$ over the whole grid range, so the small-$\zeta$ restriction of the mean-field form is removed (exact-variance note below).

\emph{Step 2 (convexity and minimizer).}
Both terms of \eqref{eq:app-window-risk} are strictly convex on $(0, 2)$, since
\[
  R''(\omega) \;=\; \frac{6\nu^2_\Sigma}{\omega^4} + \frac{4c'}{(2-\omega)^3} \;>\; 0,
\]
so $R$ has a unique minimizer.
The first-order condition is $2c'/(2-\omega)^2 = 2\nu^2_\Sigma/\omega^3$, i.e.\ $c' \omega^3 = \nu^2_\Sigma (2 - \omega)^2$; to leading order in $\omega$ the right side is $4\nu^2_\Sigma$, giving $\omega^* = (4\nu^2_\Sigma/c')^{1/3}$.
For the leading-order (cubic) model $\widetilde R(\omega) = c\,\omega + \nu^2/\omega^2$ with the critical condition $\omega^{*\,3} = 2\nu^2/c$ (i.e.\ $c = c'/2$), optimality holds by the exact factorization
\[
  \widetilde R(\omega) - \widetilde R(\omega^*)
  \;=\; \frac{\nu^2}{\omega^2 \omega^{*\,3}}\,(\omega - \omega^*)^2\,(2\omega + \omega^*) \;\ge\; 0,
\]
which is formalized (\texttt{window\_optimum} in \texttt{Ogd.lean}); substituting $c = c'/2$ reproduces the 4-form above.

\emph{Step 3 (drift inversion).}
At quasi-static equilibrium the prequential residual excess equals the full steady-state risk---lag and variance together:
\[
  \E[r_t^2] - \sigma^2
  \;=\; c'\,\frac{\omega}{2-\omega} + \frac{\nu^2_\Sigma}{\omega^2}
  \;=\; R(\omega),
\]
because the residual $r_t = y_t - x_t^\top W_{t-1}$ carries both the estimation variance and the lag of the solve.
Solving for the drift power at the current window gives
\[
  \widehat\nu^2_\Sigma
  \;=\; \max\Bigl(0,\;\bigl(\widehat E_t - \widehat c'\,\tfrac{\omega}{2-\omega}\bigr)\,\omega^2\Bigr),
  \qquad
  \widehat E_t = \widehat\sigma^2_{\mathrm{fast}} - \widehat\sigma^2_{\mathrm{slow}},
\]
where $\widehat\sigma^2_{\mathrm{fast}}, \widehat\sigma^2_{\mathrm{slow}}$ are the fast and slow residual EMAs.
One subtlety is forced by the estimator's own construction: the slow baseline is itself an average of $r_t^2$, so at stationarity $\E[\widehat\sigma^2_{\mathrm{slow}}] = \sigma^2 + \widehat c' \omega_{\mathrm{slow}}/(2 - \omega_{\mathrm{slow}})$, not the pure noise floor $\sigma^2$.
Reading the first parenthesis as the two-scale excess $\widehat E_t$ (not as $\E[r_t^2] - \sigma^2$ with a hypothetical pure noise-floor estimate) matters because, on drift onset, $\E[\widehat E_t] \approx \nu^2_\Sigma/\omega^2 + c'\omega/(2-\omega) - c'\omega_{\mathrm{slow}}/(2-\omega_{\mathrm{slow}})$, and subtracting the current window's floor isolates $\nu^2_\Sigma/\omega^2$ up to the pre-drift floor $c'\omega_{\mathrm{slow}}/(2-\omega_{\mathrm{slow}})$---negligible exactly when the pre-drift window sits at $\omega_{\min}$, the quasi-static regime of Theorem~\ref{thm:window}.
The subtraction still matters: an estimator that omits the floor term, such as $(\widehat\E[r_t^2] - \widehat\sigma^2)\omega^2$, is strictly increasing in $\omega$ (it equals $\nu^2_\Sigma + c'\omega^3/(2-\omega)$ at equilibrium) and therefore admits no positive fixed point of the estimator iteration $\omega \mapsto \omega^*(\widehat\nu^2_\Sigma)$, driving the window to its cap under sustained oscillation, whereas the variance-floor-subtracted estimator is exact at equilibrium and the iteration converges to the minimizer of Step 2; under stationarity $\E[\widehat E_t] \to 0$ and the clip pins the window at $\omega_{\min}$, the conservative direction of Step 4.

\emph{Two-scale timescale separation.}
The excess detector presupposes $h_f \ll \omega^{-1} \ll T_{\mathrm{drift}} \ll h_{\mathrm{slow}}$: the fast average settles at the post-onset level while the slow average still holds the pre-drift level, so $\widehat E_t$ isolates the drift excess.
The default $h_{\mathrm{slow}} = 20{,}000 \approx 2T$ enforces this on the grid; a drift persisting beyond $h_{\mathrm{slow}}$ would relax both averages toward the same level and close the window.

\emph{Plug-in gap in the window iteration.}
Theorem~\ref{thm:window} is stated with the exact coefficient $c' = \sigma^2 d_{\mathrm{eff}}(\rho)$, $d_{\mathrm{eff}} = \sum_i (1-s_i)^2$, whereas the runtime rule uses the mean-field plug-in $\widehat c' = \sigma^2 D(1-\bar s)^2$.
Jensen gives $\sum_i (1-s_i)^2 \ge D(1-\bar s)^2$, so the plug-in underestimates the variance coefficient; the cubed-coordinate iteration then has gain $2(c'_{\mathrm{true}}/\widehat c' - 1)$ at leading order, contracting iff $c'_{\mathrm{true}}/\widehat c' < 3/2$.
The gap is exactly the population variance of the pulls,
\[
\sum_{i=1}^D (1-s_i)^2 - D(1-\bar s)^2
= \sum_{i=1}^D (s_i-\bar s)^2 = D\operatorname{Var}(s) \;\ge\; 0,
\]
an identity formalized as \texttt{deff\_sub\_plugin} and \texttt{deff\_sub\_plugin\_nonneg} in \texttt{Core.lean} (Appendix~\ref{app:lean}); it vanishes on an isotropic spectrum and is largest when the pulls are spread out, as on a rank-deficient spectrum whose dead directions carry $s_i = 1$.
The ratio is bounded by the spread of the per-direction shrinkages and reaches $2$ in the $50\%$-dead-direction corner, so on strongly rank-deficient spectra the plug-in can drive the window to its cap.
The exact $d_{\mathrm{eff}}$ is available from the eigenvalues already computed at each refresh, at no extra matrix cost; recording the identity makes explicit that the plug-in is the conservative (over-forgetting) endpoint and that the exact coefficient is a free refinement.

\emph{Step 4 (safety asymmetry).}
Evaluating $R$ at $c\,\omega^*$ with $c > 1$ (drift overestimated, window too short) costs $O(c - 1)$ through the linear variance term, while $c < 1$ (drift underestimated, window too long) pays the quadratic lag blow-up $c^{-2}$.
As in Proposition~\ref{prop:minimax}, the conservative choice is to overestimate the adversary.
This mirrors the measured flatness of the Kalman $q$-sweep and the divergence of the leaky $\lambda$-sweep one step below its best stable window (Section~\ref{sec:experiments}).
The tradeoff itself is classical---misadjustment versus lag, from the LMS analysis of \citet{widrow1976}, whose optimum is a square-root law for random-walk drift---and the $1/3$-power law for linear drift shares the $2/3$ exponent of the dynamic-regret rate of \citet{besbes2015}. The constant is specific to the anchored accumulator, and the prequential inversion of Step 3 allows online estimation without known drift parameters.

\medskip
\noindent\textbf{Proof of Corollary~\ref{cor:mixed} (mixed-drift master law).}
For target increments $\Delta w^*_t = \nu + \vartheta_t$ with $\vartheta_t$ a zero-mean martingale difference, the exponentially weighted data path at equilibrium has tracking error $\widehat w_t - w^*_t = -\sum_{j \ge 0}(1-\omega)^j\,\Delta w^*_{t-j}$ (swap the order of summation: $\sum_{k=j}^{\infty} \omega(1-\omega)^k = (1-\omega)^j$).
The expected squared lag in prediction units is then
\[
  \mathcal{L}_{\mathrm{lag}}(\omega) = \E\Bigl\|\sum_{j \ge 0}(1-\omega)^j\,\Sigma^{1/2}(\nu + \vartheta_{t-j})\Bigr\|^2
  = \frac{\nu^2_\Sigma}{\omega^2} + \frac{q_\Sigma}{\omega(2-\omega)},
\]
where the cross term vanishes (martingale difference) and the geometric sums evaluate in closed form ($\sum_j (1-\omega)^j = 1/\omega$, $\sum_j (1-\omega)^{2j} = 1/(\omega(2-\omega))$).
The master risk $R(\omega) = c' \omega/(2-\omega) + q_\Sigma/(\omega(2-\omega)) + \nu^2_\Sigma/\omega^2$ specializes to: (i) the coherent regime ($q_\Sigma = 0$), giving Theorem~\ref{thm:window}'s $\omega^* = (4\nu^2_\Sigma/c')^{1/3}$; (ii) the random-walk regime ($\nu_\Sigma = 0$), giving $\omega^* = (q_\Sigma/c')^{1/2} + O(q_\Sigma/c')$, the classical $1/2$-power law.
The runtime estimator inverts the total lag without knowing the regime: $\widehat\nu^2_\Sigma = \max\bigl(0,\,(\widehat E_t - \hat c' \omega/(2-\omega))\,\omega^2\bigr)$ with $\widehat E_t$ the two-scale residual excess of Step 3 (Appendix~\ref{app:proofs}) estimates the drift excess above the variance floor regardless of whether it originates from coherent or stochastic drift.

\emph{Contraction in the random-walk regime.}
When $\nu_\Sigma = 0$ and $q_\Sigma > 0$, the excess above the variance floor is $q_\Sigma/(\omega(2-\omega))$, so the estimator returns $\widehat\nu^2_\Sigma(\omega) = q_\Sigma\,\omega/(2-\omega)$ and the cubic map $T_{\mathrm{RW}}(\omega) = (4\widehat\nu^2_\Sigma(\omega)/c')^{1/3}$ satisfies
\[
  \bigl(T_{\mathrm{RW}}(\omega)\bigr)^3 = \frac{4q_\Sigma\,\omega}{c'\,(2-\omega)}.
\]
In logarithmic coordinates $\frac{d\ln T_{\mathrm{RW}}}{d\ln \omega} = \frac13\bigl(1 + \frac{\omega}{2-\omega}\bigr) = \frac{2}{3(2-\omega)} \le \frac23 < 1$, so $T_{\mathrm{RW}}$ is a strict contraction.
Its unique positive fixed point solves $\omega^2(2-\omega) = 4q_\Sigma/c'$, hence $\omega_{\mathrm{fp}} = \sqrt{2q_\Sigma/c'}\,(1+O(\omega_{\mathrm{fp}})) = \sqrt2\,\omega^*_{\mathrm{RW}}$.
At that point the random-walk risk $R_{\mathrm{RW}}(\omega) = c'\omega/(2-\omega) + q_\Sigma/(\omega(2-\omega))$ equals $\frac{3}{2\sqrt2} \approx 1.061$ times its value at $\omega^*_{\mathrm{RW}}$ in the limit $q_\Sigma/c'\to 0$ (finite ratios give larger factors, and the fixed point exists for $q_\Sigma/c' \le 8/27$), so the same $1/3$-power update achieves the exact linear-drift optimum and at most $6.1\%$ asymptotic overhead under random-walk drift without knowing the regime.
\end{proof}

\medskip
\noindent\emph{Exact finite-window variance term.}
The mean-field term $c'\,\omega/(2-\omega)$ is the $\zeta \to 0$ limit of the exact random-design variance $\sigma^2\Upsilon(\omega;\rho)/(1-\Upsilon(\omega;\rho))$ of Theorem~\ref{thm:rmt}(i): the expansion of Appendix~\ref{app:rmt-variance} gives $\sigma^2\Upsilon/(1-\Upsilon) \to \tfrac12\sigma^2\omega\,d_{\mathrm{eff}}(\rho)$.
Using the exact form in the cycle-averaged risk,
\[
R(\omega,\rho) \;=\; \frac{\sigma^2\,\Upsilon(\omega;\rho)}{1-\Upsilon(\omega;\rho)} + \frac{\nu^2_\Sigma}{\omega^2},
\]
removes the small-$\zeta$ restriction: over the whole grid range $\zeta \in [0.26, 6.4]$ the exact form matches the Monte Carlo stationary risk to within $5\%$ (median $2.4\%$), whereas the mean-field term deviates by up to $100\%$ at $\zeta = 6.4$ (\texttt{diagnostics.py} \texttt{--diag exactrisk}; \texttt{results/exact\_risk.csv}).
The runtime keeps the mean-field plug-in $\hat c'_t$, whose Jensen gap is the conservative direction of Appendix~\ref{app:proofs}; the exact form is the analysis-level statement of the window law, available at every window from the closing eigenvalues, not a runtime change.

\subsection{Proof of Proposition~\ref{prop:conflict}(i) (oracle dominance)}
\label{app:proof-dominance}



\begin{proof}[Proof of Proposition~\ref{prop:conflict}(i)]
We first record the risk-interpolation step used throughout.
By the pull identity (Equation~\eqref{eq:pull}) with $\gamma = 1$,
\[
  W - w \;=\; (1 - s)\,(\widehat w - w) + s\,\delta .
\]
Squaring and taking expectations, the cross term is $2(1-s)s\,\delta\,\E[\widehat w - w]$, which vanishes conditioned on $\mathcal F_{t_0}$ when the anchor $G$ is frozen (drawn before the post-anchor data); at finite horizons with a refreshed anchor, Theorem~\ref{thm:cadence} bounds the induced correlation.
\begin{equation}
\label{eq:app-risk-interp}
  \E\bigl[(W - w)^2\bigr] \;=\; (1 - s)^2\,v + s^2\,\delta^2 .
\end{equation}
Now set $\delta = 0$: the second term disappears and \eqref{eq:app-risk-interp} gives $\E[(W_i - w_i)^2] = (1 - s_i)^2 v_i < v_i$ whenever $0 < s_i < 1$ and $v_i > 0$.
Since the per-direction prediction risk contribution is $\mu_i\,\E[(W_i - w_i)^2]$, summing gives the display.
The improvement factor per direction is $\bigl(\mu_i/(\mu_i + \rho)\bigr)^2$, which is smallest (largest gain) when $\mu_i \ll \rho$.
The interpolation identity \eqref{eq:app-risk-interp} and both dominance statements are formalized (\texttt{convex\_interp\_sq}, \texttt{dominance}, \texttt{dominance\_strict} in \texttt{Core.lean}).
\end{proof}

\begin{remark}[The $\gamma = 1$ design principle]
With $\gamma < 1$ the anchored estimate is pulled toward $\gamma g$ rather than $g$ (Equation~\eqref{eq:pull}), introducing the shrinkage bias $s(1-\gamma)|w|$ per direction, which dominates the variance benefit whenever the prior is informative.
All methods therefore use $\gamma = 1$; the ablation measures the cost of $\gamma = 0.5$ directly ($+560\%$, Online Appendix~\ref{app:protocol}).
\end{remark}

\subsection{Proof of Proposition~\ref{prop:conflict}(ii, first half)}
\label{app:proof-imperfect}



\begin{proof}[Proof of Proposition~\ref{prop:conflict}(ii)]
(a) For real $u, d$ and $s \in [0,1]$, expanding the square gives
\[
  \bigl((1-s)u + s d\bigr)^2
  \;=\; (1-s)u^2 + s d^2 - s(1-s)(u - d)^2
  \;\le\; (1-s)u^2 + s d^2 ,
\]
and taking expectations with $u = \widehat w - w$, $d = \delta$ applied to the error decomposition of Equation~\eqref{eq:pull} yields the first inequality; the rest is elementary.
The exact deficit $s(1-s)(u-d)^2$ is formalized as \texttt{convex\_interp\_sq}.
(b) Completing the square,
\[
  \frac{vc}{v+c} + (v+c)\Bigl(s - \frac{v}{v+c}\Bigr)^{2}
  \;=\; (1-s)^2 v + s^2 c
\]
by pure algebra, which is the stated identity.
The excess $(v+c)(s - s^*)^2$ is nonnegative with equality iff $s = s^*$, giving uniqueness and the cost bound.
The identity and the optimal pull are formalized (\texttt{quadratic\_calibration}, \texttt{optimal\_pull\_min}, \texttt{optimal\_pull\_attained} in \texttt{Core.lean}).
\end{proof}

\subsection{Proof of Proposition~\ref{prop:conflict}(ii, second half)}
\label{app:proof-sign}



\begin{proof}[Proof of Proposition~\ref{prop:conflict}(ii), strength form]
(i) By Proposition~\ref{prop:conflict}(ii) the optimal pull in direction $i$ is $s^* = v_i/(v_i + \delta_i^2)$, and $\rho \mapsto s = \rho/(\mu_i + \rho)$ is a strictly increasing bijection from $[0,\infty)$ onto $[0,1)$ with inverse $\rho = \mu_i s/(1-s)$.
Substituting $s^*$ gives $\rho^*_i = \mu_i v_i/\delta_i^2 = \sigma^2/(n\delta_i^2)$.
(ii) With a single pull $s$ the total risk is $R(s) = D(1-s)^2 v + s^2 C$ with $v = \sigma^2/(n\mu)$; minimizing as in Proposition~\ref{prop:conflict}(ii) gives $s^* = Dv/(Dv + C)$, hence $\rho^* = \mu D v/C = \sigma^2 D/(nC)$.
(iii) Immediate by differentiating the closed forms of (i)--(ii).
As $n \to \infty$ at fixed conflict the optimal strength vanishes, and under stationarity the measured conflict decays to its noise floor while the optimal strength recovers: a rising drift estimate lowers the anchor strength toward the floor, and a stationary stream lets it return.
\end{proof}

\begin{remark}[Runtime rule: prediction units]
The runtime rule is $\rho^* = \bar\mu_t\,\widehat V^\Sigma_t/\widehat C^\Sigma_t$ with $\widehat V^\Sigma_t = \widehat\sigma^2_t\,d_{\mathrm{act}}/n_t$ and $\widehat C^\Sigma_t = \max\{C^\Sigma_{\mathrm{ref}}, C^\Sigma_{\mathrm{resid}}, C^\Sigma_{\mathrm{floor}}\}$; per direction it reduces to the closed form of Proposition~\ref{prop:conflict}(ii), $\mu_i \hat v_i/\hat c_i = \sigma^2/(n\delta_i^2)$.
Every branch is in prediction-risk units $[y^2]$: the reference branch is $C^\Sigma_{\mathrm{ref}} = (G_{\mathrm{ref}} - \widehat w)^\top \bar A_t (G_{\mathrm{ref}} - \widehat w) - \widehat V^\Sigma$, the residual branch is $C^\Sigma_{\mathrm{resid}} = E_t$, and the floor is $C^\Sigma_{\mathrm{floor}} = \tfrac12 \widehat\sigma^2 d_{\mathrm{act}}/T_w$.
The naive conflict $\|G - \widehat w\|^2$ has expectation $C + V$ under the null, so the estimate is debiased by subtracting $\widehat V^\Sigma$; because a refreshed anchor co-moves with the current solve, conflict is measured against the slow reference $G_{\mathrm{ref}}$, while the two-scale residual detector (Proposition~\ref{prop:falsealarm} its calibration model) supplies fast disengagement.
Under deep stationarity $\widehat C^\Sigma \to 0$ would drive $\rho^* \to \infty$; the floor and a hard cap remove the singularity.
The \emph{legacy parameter-unit} form uses $\widehat V = \widehat\sigma^2\,\mathrm{tr}(\bar A^{-1})/n$ and $C_{\mathrm{resid}} = E_t\,\mathrm{tr}(\bar A^{-1})/(n\bar\mu)$; the trace scales with the effective condition number rather than with $d_{\mathrm{act}}$, so it inflates on ill-conditioned spectra and is reported only as the ablation of Online Appendix~\ref{app:protocol}.
Overestimating the conflict is safe by Proposition~\ref{prop:minimax}.
\end{remark}

\medskip
\noindent\emph{Scalar-$\rho$ prediction risk on a non-isotropic spectrum.}
Proposition~\ref{prop:conflict}(ii) gives the per-direction strength, while Algorithm~\ref{alg:spar-pull} runs a single scalar $\rho$; the scalar target is obtained by differentiating the prediction risk directly.
With $s_i = \rho/(\mu_i+\rho)$ and $v_i = \sigma^2/(n\mu_i)$,
\[
  R_{\mathrm{pred}}(\rho)
  = \sum_{i=1}^D \mu_i\bigl[(1-s_i)^2 v_i + s_i^2\delta_i^2\bigr]
  = \sum_{i=1}^D \frac{\mu_i^2(\sigma^2/n) + \rho^2\,\mu_i\delta_i^2}{(\mu_i+\rho)^2},
\]
so
\[
\begin{aligned}
  \frac{d R_{\mathrm{pred}}}{d\rho}
  &= 2\sum_{i=1}^D \frac{\mu_i^2}{(\mu_i+\rho)^3}
    \Bigl(\rho\,\delta_i^2 - \frac{\sigma^2}{n}\Bigr) = 0 \\
  \Longleftrightarrow\quad
  \rho^* &= \frac{(\sigma^2/n)\sum_i \psi_i(\rho^*)}{\sum_i \psi_i(\rho^*)\,\delta_i^2},
  \qquad
  \psi_i(\rho) := \frac{\mu_i^2}{(\mu_i+\rho)^3},
\end{aligned}
\]
which weights directions near the crossover $\mu_i \approx 2\rho$; this is the exact scalar optimum, and it is generally not $\bar\mu\widehat V^\Sigma/\widehat C^\Sigma$.
Replacing the exact sensitivity $\psi_i(\rho)$ by the excitation-proportional surrogate $\psi_i \propto \mu_i$ gives the plug-in approximation
$\rho_{\mathrm{plug}} = (\sigma^2/n)\sum_i \mu_i \big/ \sum_i \mu_i\,\delta_i^2$, which in prediction units with $\widehat V^\Sigma = \sigma^2 d_{\mathrm{act}}/n$, $C^\Sigma = \sum_i \mu_i\delta_i^2$, and $\bar\mu = \mathrm{tr}(\bar A)/d_{\mathrm{act}}$ is the implemented rule $\rho_{\mathrm{plug}} = \bar\mu\,\widehat V^\Sigma/\widehat C^\Sigma$.
The surrogate matches the exact sensitivity where it matters: $\psi_i(\rho) = O(\mu_i^2/\rho^3)$ as $\mu_i \to 0$, so both vanish on weakly excited and dead directions and the conflict is carried by the well-excited modes, whereas uniform weights would give a dead direction the same weight as an excited one.
The two rules coincide on isotropic spectra, where the direction weights are equal.

\medskip
\noindent\emph{Crossover horizon for the plug-in estimator.}
By the calibration identity, a static pull $s_0$ pays $(v + c)\,(s_0 - s^*)^2$ over the calibrated optimum, while the plug-in estimate $\hat s$ pays $(v + c)\,(\hat s - s^*)^2$.
If the two-scale noise and conflict estimates are consistent with $\lvert\hat s_t - s^*\rvert \le K\,t^{-1/2}$, the plug-in beats any fixed pull once
\[
  t \;>\; T_{\mathrm{crit}} \;\asymp\; \frac{K^2}{(s_0 - s^*)^2},
\]
the square of the calibration constant over the squared pull mismatch (the risk scale $v+c$ cancels on both sides).
Short streams therefore favor fixed anchoring, and the BCI IV streams ($T \approx 300$ trials) sit on the short side; on the grid the fixed-$\rho$ ablation is $2.2\%$ better at $T = 10^4$, so our runs do not resolve the crossover. The threshold scales quadratically in the prior mismatch $s_0 - s^*$; the plug-in-versus-fixed-shrinkage tradeoff is classical in empirical-Bayes shrinkage \citep{ledoitwolf2004}.

\subsection{Proof of Proposition~\ref{prop:minimax} (minimax safety)}
\label{app:proof-minimax}



\begin{proof}[Proof of Proposition~\ref{prop:minimax}]
For fixed $C$, $R$ is strictly convex in $s$ with unique minimizer $s^*(C) = Dv/(Dv + C)$ (Proposition~\ref{prop:conflict}(ii) applied to the lumped problem; the interpolation identity is Proposition~\ref{prop:conflict}(i)), and $\rho \mapsto s(\rho)$ is a bijection, so $\min_\rho R(\rho, C) = DvC/(Dv + C)$, attained uniquely at $\rho^*(C) = \mu Dv/C$.
This value is strictly increasing in $C$, so the maximin is attained at $C_0$ and equals $DvC_0/(Dv + C_0)$.
For the minimax: for fixed $\rho$ with $s(\rho) > 0$, $R(\rho, C)$ is increasing in $C$, so $\max_{C \le C_0} R(\rho, C) = R(\rho, C_0)$, and minimizing over $\rho$ gives the same value, attained at $\rho^*(C_0)$.
The inequality $DvC/(Dv + C) \le (1 - s)^2 Dv + s^2 C$ for all $s$, which is the computational core of the equality (together with the monotonicity of the value in $C$), is formalized (\texttt{minimax\_lower} in \texttt{Core.lean}).
\end{proof}



\begin{proof}[Proof of Proposition~\ref{prop:minimax}, safety clause]
(i) is Proposition~\ref{prop:conflict}(ii).
(ii) Expand: $R(s) - v = (1-s)^2 v + s^2 C - v = s\bigl(s(v + C) - 2v\bigr)$, which is at most $0$ exactly when $s \le 2v/(v+C) = 2s^*$.
Since $\widehat C \ge C$ implies $\widehat s = v/(v + \widehat C) \le v/(v + C) = s^* < 2s^*$, we get $R(\widehat s) \le v$.
(iii) As $\widehat C \to \infty$, $\rho^* = \widehat V/\widehat C \to 0$ and $s_i = \rho/(\mu_i + \rho) \to 0$ for every direction with $\mu_i > 0$; directions with $\mu_i = 0$ contribute $\mu_i \delta_i^2 = 0$ to the prediction risk.
\end{proof}

\subsection{Proof of Proposition~\ref{prop:falsealarm} (false-alarm calibration)}
\label{app:proof-falseesc}



\begin{proof}[Proof of Proposition~\ref{prop:falsealarm}]
The fast two-scale excess is the causal moving average $Z_t = \sum_{k\ge 0} w_k z_{t-k}$ of a centered sub-exponential martingale-difference sequence, with geometric weights $w_k = \eta_f(1-\eta_f)^k$; its weight energy is the geometric sum
\[
\sum_{k\ge 0} w_k^2 = \frac{\eta_f^2}{1-(1-\eta_f)^2} = \frac{\eta_f}{2-\eta_f} = \frac{1}{m_{\mathrm{eff}}},
\qquad m_{\mathrm{eff}} = \frac{2-\eta_f}{\eta_f},
\]
and $\max_k w_k = \eta_f$.
Freedman's inequality for sub-exponential martingale differences \citep{freedman1975} bounds the weighted sum at each time with predictable variance $V = \nu_z^2/m_{\mathrm{eff}}$ and scale $b = \eta_f b_z$:
\[
\Pr[Z_t > \tau] \;\le\; \exp\Bigl(-\frac{\tau^2}{2\,(V + b\,\tau)}\Bigr),
\]
so the threshold $\tau(\delta) = bL + \sqrt{b^2L^2 + 2VL}$ with $L = \log(T/\delta)$ solves $\tau^2 = 2(V + b\tau)L$ and gives $\Pr[Z_t > \tau(\delta)] \le e^{-L} = \delta/T$.
A union bound over the $T$ time points upgrades this to the sliding (time-uniform) statement $\Pr[\exists\, t \le T: Z_t > \tau(\delta)] \le T \cdot \delta/T = \delta$.
The statistic is the implemented one, with effective sample size $m_{\mathrm{eff}} = (2-\eta_f)/\eta_f \to 1/\eta_f$ for small $\eta_f$, the runtime constant; no disjoint-window surrogate is used.
The damage of a false escalation is bounded by the cost of running at reduced anchor strength for the escalation holdout: at most $O(v)$ per step for $O(1/\eta_f)$ steps, so with $\delta = 1/T^2$ the total expected false-escalation damage is $o(1)$.
The alternative (drift) is detected with lag $O(h_f + 1/\eta_f)$, where $h_f$ is the fast moving-average horizon.
\end{proof}

\medskip
\noindent\emph{Stationary pinning (effective threshold).}
The variance-floor subtraction is itself a strictly positive threshold: the drift-energy estimator is $\widehat\nu^2_\Sigma=\max\bigl(0,(E_t-\tau_\omega)\omega^2\bigr)$ with $\tau_\omega=\hat c'\,\omega_{\min}/(2-\omega_{\min})$, so $E_t\le\tau_\omega$ implies $\widehat\nu^2_\Sigma=0$ and $\omega_t=\omega_{\min}$; the conflict floor $\tau_\rho=C_{\mathrm{floor}}$ likewise prevents the residual branch from driving the strength below its stationary floor.
Applying the bound of Proposition~\ref{prop:falsealarm} at $\tau=\min(\tau_\omega,\tau_\rho)>0$ gives $\Pr[\omega_t>\omega_{\min}]\le\delta$ for the sliding test under stationarity, upgrading Theorem~\ref{thm:window}'s in-expectation statement to a high-probability pinning guarantee; the bound applies to the implemented EMA directly (Proposition~\ref{prop:falsealarm}).

\begin{remark}[Time uniformity]
The union bound over time points makes the guarantee uniform over the whole stream, so no disjoint-window or e-process surrogate is needed; the cost is the $\log T$ factor in $\tau(\delta)$.
The measured rate is at or below the nominal level at $\delta = 0.10$ and $0.05$, with a slight exceedance at $\delta = 0.01$ on ill-conditioned spectra (Online Appendix~\ref{app:protocol}).
\end{remark}

\subsection{Proof of Proposition~\ref{prop:ocowrap}(i) (projected gradient)}
\label{app:proof-composition}



\begin{proof}[Proof of Proposition~\ref{prop:ocowrap}(i)]
\emph{Step 1 (convexity and subgradient).}
Writing $z_t = U_t^\top x_t$, $w^e_t = U_t^\top \widehat w_t$, and $g^e_t = U_t^\top G_t$, the prediction is $a_t(s) := x_t^\top W_t(s) - y_t = z_t^\top[(1-s)\odot w^e_t + s\odot g^e_t] - y_t$, affine in $s$ on the box; hence $\ell_t = a_t(s)^2$ is convex, with
\[
  \nabla_s \ell_t(s) \;=\; 2\,a_t(s)\,\bigl(z_{t,i}\,(g^e_{t,i} - w^e_{t,i})\bigr)_i .
\]
The exact subgradient identity for the segment family, $(a + b w)^2 + 2b(a + bw)(s - w) \le (a + b s)^2$, is formalized (\texttt{sq\_subgradient} in \texttt{Ogd.lean}), which is the convexity step in the Lean assembly below.

\emph{Step 2 (gradient bound).}
Since $s_i \in [0,1]$ each eigen-coordinate obeys $|W^e_{t,i}| \le \max\{|w^e_{t,i}|, |g^e_{t,i}|\}$, and the Euclidean norm is invariant under $U_t$, so
\[
  \|W_t(s)\|_2^2 \;=\; \|W^e_t(s)\|_2^2 \;\le\; \|\widehat w_t\|_2^2 + \|G_t\|_2^2 \;\le\; 2W_{\max}^2,
\]
and therefore
\[
  |a_t(s)| \;\le\; \|x_t\|_2\, \|W_t(s)\|_2 + |y_t| \;\le\; \sqrt 2\, X W_{\max} + Y .
\]
Using $\|z_t\|_\infty \le \|z_t\|_2 = \|x_t\|_2 \le X$ and $\|g^e_t - w^e_t\|_2 = \|G_t - \widehat w_t\|_2 \le 2W_{\max}$,
\[
  \bigl\|\bigl(z_{t,i}(g^e_{t,i} - w^e_{t,i})\bigr)_i\bigr\|_2
  \;\le\; \|z_t\|_\infty\, \|g^e_t - w^e_t\|_2
  \;\le\; 2\,X W_{\max},
\]
and combining the two factors gives $\|\nabla_s \ell_t\|_2 \le G_b$ with $G_b = 4XW_{\max}(\sqrt 2 XW_{\max} + Y)$.

\emph{Step 3 (projected descent).}
The box $[0,1]^D$ has $\ell_2$-diameter $\sqrt D$, and for any comparison point $s^*$ in the box, $\|s_0 - s^*\|_2 \le \sqrt D$.
Online projected gradient descent on convex losses against an adaptive adversary incurs regret at most
\[
  \frac{\|s_0 - s^*\|_2^2}{2\iota} + \frac{\iota\, T\, G_b^2}{2}
  \;\le\; G_b\sqrt{DT}
  \;\le\; G_b\sqrt{2DT}
\]
with the stated step size \citep{zinkevich2003}; no probabilistic assumption is used.
The comparison class contains the pure-data profile $\mathbf 0$, the pure-anchor profile $\mathbf 1$, and every fixed scalar-$\rho$ profile $s_\rho$, since all lie in the box.
The bound is external regret against the realized loss sequence: the anchor's state makes a comparator's counterfactual anchor trajectory differ from the realized one, and policy regret is not claimed.
\end{proof}

\begin{remark}[Formalized chain]
The descent lemma and its assembly are mechanized in \texttt{Ogd.lean}:
\texttt{sq\_subgradient} (the convexity step),
\texttt{projBox\_dist\_sq\_le} (box-projection contraction),
\texttt{ogd\_regret\_box} (the Zinkevich step: regret at most $(w_0 - s^*)^2/(2\eta) + \eta T G^2/2$ against any fixed comparison point, with adaptive losses and no probability),
and \texttt{ogd\_regret\_segment} (the composition assembled for the scalar-pull family).
The $D$-dimensional status is described in Appendix~\ref{app:lean}.
\end{remark}

\subsection{Proof of Proposition~\ref{prop:ocowrap}(ii) (Hedge)}
\label{app:proof-hedge}



\begin{proof}[Proof of Proposition~\ref{prop:ocowrap}(ii)]
\emph{Step 1 (quadratic exponential bound).}
For $x \ge 0$,
\[
  e^{-x} \;\le\; \frac{1}{1+x} \;\le\; 1 - x + x^2,
\]
the first inequality by $e^x \ge 1 + x$, the second by the exact identity $(1 - x + x^2)(1 + x) = 1 + x^3 \ge 1$.
This step is formalized (\texttt{exp\_neg\_le\_one\_sub\_add\_sq} in \texttt{Hedge.lean}).

\emph{Step 2 (potential ratio).}
By the closed form of the weights,
\[
  Z_{t+1}
  \;=\; \sum_i e^{-\eta(L_{i,t} + \ell_{i,t})}
  \;=\; Z_t \sum_i \pi_{i,t}\, e^{-\eta \ell_{i,t}} ,
\]
which is formalized as \texttt{Zfun\_step}.
Applying Step 1 termwise with $\ell_{i,t} \le M$ (so $\ell_{i,t}^2 \le M\,\ell_{i,t}$), and using $\sum_i \pi_{i,t} = 1$ and $1 + u \le e^u$,
\[
  \frac{Z_{t+1}}{Z_t}
  \;\le\; \sum_i \pi_{i,t}\bigl(1 - \eta \ell_{i,t} + \eta^2 \ell_{i,t}^2\bigr)
  \;\le\; 1 - \eta(1 - \eta M)\,m_t
  \;\le\; e^{-\eta(1 - \eta M)\,m_t}
\]
(\texttt{Zfun\_step\_bound}).

\emph{Step 3 (telescope).}
Multiplying over $t < T$ and using $Z_0 = K$ (\texttt{Zfun\_prod\_bound}) gives
$Z_T \le K\, e^{-\eta(1 - \eta M)\sum_t m_t}$, while the $j$-th term of the defining sum gives $Z_T \ge e^{-\eta L_{j,T}}$.
Taking logarithms,
\[
  e^{-\eta L_{j,T}}
  \;\le\; K\, e^{-\eta(1 - \eta M)\sum_t m_t}
  \qquad\Longrightarrow\qquad
  (1 - \eta M)\sum_{t < T} m_t \;\le\; L_{j,T} + \frac{\ln K}{\eta},
\]
which is the stated bound since $1 - \eta M > 0$ (\texttt{hedge\_regret}).
The mixture-convexity step $(\sum_i \pi_i u_i)^2 \le \sum_i \pi_i u_i^2$ is the interpolation bound of Proposition~\ref{prop:conflict}(i), formalized as \texttt{convex\_interp\_sq\_le}, so the realized mixture loss is bounded per step; the anytime form (the implemented schedule $\varsigma_t = \sqrt{\ln K / S_t}$) is the standard AdaHedge substitution \citep{cesabianchi2006}.
\end{proof}

\medskip
\noindent\emph{Best of both worlds.}
Because \texttt{hedge\_regret} holds for every comparator $j$ in the pool, the frozen anchor and the best adapted window can be chosen simultaneously:
\[
  \sum_{t<T} \ell_t(\mathrm{mix}) \;\le\; \frac{1}{1-\eta M}\Bigl(\min_j L_{j,T} + \frac{\ln K}{\eta}\Bigr),
\]
with the pool containing both: the mixture is never worse than the warm model by more than the regret term, and never worse than the best adapted path by the same term. The SPAR-EG$^+$ pool includes both the frozen anchor and the plug-in method.
\citet{evendar2008} quantify the regret-to-best versus regret-to-average tradeoff.

\subsection{Proof of Proposition~\ref{prop:windup} (conditioning bound)}
\label{app:proof-props}



\begin{proof}[Proof of Proposition~\ref{prop:windup}]
The eigenvalues of $A + n\rho I$ lie in the interval $[\,n\rho,\ n(\mu_{\max} + \rho)\,]$, whose ratio is at most $1 + \mu_{\max}/\rho$.
\end{proof}

The accumulator of Section~\ref{sec:theory} never enters the covariance form: the anchored solve is a batch information-form solve, so Proposition~\ref{prop:windup} holds at every window, including windows where the covariance recursion of leaky RLS loses rank and diverges \citep{sayed2008}.
The contrast is measured in Section~\ref{sec:experiments}: the leaky covariance recursion diverges in $135/135$ main-grid and $20/20$ stress-suite runs at window $100 < D = 128$, while the batch information-form solve of the same objective is stable at every $\lambda$, and no SPAR method diverged in any configuration.

\medskip
\noindent\emph{Eigengap context (Davis--Kahan).}
Let $U, U'$ be the eigenbases of $\bar A$ and $\bar A' = \bar A + E$ with $\|E\|_{\mathrm{op}} \le \epsilon$, and let $\delta$ be the minimum gap between eigenvalues of $\bar A$ involved in a per-direction decision.
Then
\[
  \bigl\|(I - U'U'^\top)\,U_i\bigr\| \;\le\; \frac{\epsilon}{\delta}
\]
by the Davis--Kahan $\sin\Theta$ theorem in the two-subspace form \citep{davis1970,wedin1972}, so the per-direction refresh with a stale basis errs by at most $O(\epsilon/\delta)\,\|W - G\|$ per refresh cycle.
A refresh cycle performed in a basis rotated by more than the pull displacement applies the intended per-direction weights to the wrong coordinates, at a cost proportional to the rotation angle times $\|W - G\|$.
Three regimes follow: with static, well-separated spectra (the ill-conditioned grid) the bound is small and the per-direction estimator is stable and wins pre-drift; with flat spectra ($\delta \to 0$) the eigenbasis is essentially arbitrary and per-direction state is noise; with rotating spectra (covariance drift) the basis drifts within the estimator's memory and per-direction state scrambles.
The scalar-strength method aggregates over directions and is invariant to basis rotation, which is why it is the robust default.

\subsection{Spectral retention (Theorem~\ref{thm:rmt})}
\label{app:rmt}

This subsection separates the deterministic-equivalent proof from its numerical validation.
Throughout, $\zeta = D\omega$ is the aspect ratio held fixed in the proportional limit, $H$ is the population spectral measure of $\Sigma$, and $\rho_{\mathrm{eff}}$ is the finite-window renormalized anchor strength.
The risk coefficient of Theorem~\ref{thm:window} is written $c' = \sigma^2 d_{\mathrm{eff}}(\rho)$ throughout; it is distinct from $\zeta$ and from the cost branches $C^\Sigma_{\mathrm{ref}}, C^\Sigma_{\mathrm{resid}}, C^\Sigma_{\mathrm{floor}}$ of Algorithm~\ref{alg:spar-pull}.
The scalar geometric weights are $\omega_s = \omega(1-\omega)^s$; $W$, $w^*$, and $\widehat w$ denote the anchored solve, the target parameter, and the data solve as in the main text.

\subsubsection{The deterministic-equivalent system}
\label{app:rmt-system}

For the exponentially weighted accumulator with geometric weights $\omega(1-\omega)^s$ and features $x_s = \Sigma^{1/2} z_s$ ($z_s$ sub-Gaussian with independent entries), the separable-covariance deterministic equivalent \citep{silverstein1995,rubio2011,hachem2007} gives, for the resolvent $R(\rho) = (\bar A + \rho I)^{-1}$ and $T(-\rho) = (\kappa_1 \Sigma + \rho I)^{-1}$, the two-equation fixed-point system
\begin{equation}
\label{eq:rmt-system}
\theta_1 = \int \frac{\tau}{\kappa_1 \tau + \rho}\, dH(\tau),
\qquad
\kappa_1 = \sum_{s \ge 0} \frac{\omega_s}{1 + D \omega_s \theta_1}.
\end{equation}
In the streaming limit ($D \to \infty$, $\omega \to 0$, $D \omega = \zeta$ fixed) the weight sum integrates in closed form.
Substituting $y = (1-\omega)^s$ and using $\sum_s \omega (1-\omega)^s f((1-\omega)^s) \to \int_0^1 f(y)\,dy$,
\begin{equation}
\label{eq:rmt-kappa}
\kappa_1(\theta_1)
= \int_0^1 \frac{dy}{1 + \zeta \theta_1 y}
= \frac{\ln(1 + \zeta \theta_1)}{\zeta \theta_1}.
\end{equation}
The hypotheses used here are the ones stated in Theorem~\ref{thm:rmt}: independent sub-Gaussian entries, a population spectrum whose nonzero part is bounded away from zero ($\tau_{\min}^+ > 0$), and the proportional regime.
The atom at zero is allowed: Corollary~\ref{cor:rmt-active} below shows that exact null directions drop out of the prediction risk, so the system collapses to the active subspace verbatim.

\subsubsection{The collapse lemma}
\label{app:rmt-collapse}

\begin{lemma}[Collapse]
\label{lem:rmt-collapse}
Let $q = 1 - \bar s$ be the average data retention. Under \eqref{eq:rmt-system}--\eqref{eq:rmt-kappa}, $\kappa_1 \theta_1 + \bar s = 1$, and the system collapses to the scalar fixed point
\begin{equation}
\label{eq:rmt-collapse}
q = \int \frac{\tau}{\tau + \rho_{\mathrm{eff}}}\, dH(\tau),
\qquad
\rho_{\mathrm{eff}} = \rho\,\frac{e^{\zeta q} - 1}{\zeta q}.
\end{equation}
\end{lemma}

\begin{proof}
Per population eigenvalue the retention and the pull sum to one, so $\kappa_1 \theta_1 + \bar s = 1$ and $q = \kappa_1 \theta_1$.
Substituting \eqref{eq:rmt-kappa} gives $q = \ln(1 + \zeta \theta_1)/\zeta$, hence $e^{\zeta q} = 1 + \zeta \theta_1$.
Therefore $\kappa_1 = \zeta q/(e^{\zeta q} - 1)$ and
$\rho_{\mathrm{eff}} = \rho/\kappa_1 = \rho (e^{\zeta q} - 1)/(\zeta q)$.
The collapse identities are formalized as \texttt{de\_log\_identity} and \texttt{de\_collapse} in \texttt{Core.lean}; the probabilistic convergence of the deterministic equivalent remains pen-and-paper.
Existence and uniqueness in $(0,1)$: the right side of \eqref{eq:rmt-collapse} is continuous and strictly decreasing in $q$ (the growth factor $(e^{\zeta q}-1)/(\zeta q) \ge 1$ is strictly increasing, and each term $\tau/(\tau + \rho_{\mathrm{eff}})$ decreases in $\rho_{\mathrm{eff}}$), positive at $q \downarrow 0$ because $\tau_{\min} > 0$, and strictly below $1$ at $q = 1$.
\end{proof}

\subsubsection{Variance functional (Theorem~\ref{thm:rmt}(i))}
\label{app:rmt-variance}

With an oracle anchor ($\delta = 0$), the one-step excess prediction risk is its variance component
\begin{equation}
\label{eq:rmt-variance}
\sigma^2\,\E\!\left[\operatorname{tr}\!\left(\Sigma R\, \bar A^{(2)} R\right)\right],
\qquad
\bar A^{(2)} = \sum_s \omega_s^2 x_s x_s^\top .
\end{equation}
Perturb the weights by $\omega_s(\varepsilon) = \omega_s(1 + \varepsilon \omega_s)$, so that $\omega_s'(0) = \omega_s^2$ and $\bar A^{(2)}$ is the first-order term.
Then
\begin{equation}
\label{eq:rmt-derivative}
\E\!\left[\operatorname{tr}\!\left(\Sigma R\, \bar A^{(2)} R\right)\right]
= -\left.\frac{d}{d\varepsilon}\E\operatorname{tr}\!\left(\Sigma R_\varepsilon\right)\right|_{\varepsilon=0}
= -D\,\dot\theta_1,
\end{equation}
where the dot is $d/d\varepsilon$ at $\varepsilon = 0$.
This is the step that produces the amplification denominator.

Differentiate \eqref{eq:rmt-system} in $\varepsilon$.
Writing $J = \int \bigl(\tau/(\tau + \rho_{\mathrm{eff}})\bigr)^2 dH(\tau)$, the first equation gives
\begin{equation}
\label{eq:rmt-theta-dot}
\dot\theta_1 = -\frac{\dot\kappa_1}{\kappa_1^2}\,J,
\end{equation}
because $\kappa_1\tau + \rho = \kappa_1(\tau + \rho_{\mathrm{eff}})$.
For the second equation, the derivative of $\omega_s(\varepsilon)$ and the derivative of the denominator combine:
\begin{equation}
\label{eq:rmt-kappa-dot}
\dot\kappa_1
= \sum_s \left[
\frac{\omega_s^2}{1 + D \omega_s \theta_1}
- \frac{\omega_s \bigl(D \omega_s^2 \theta_1 + D \omega_s \dot\theta_1\bigr)}{(1 + D \omega_s \theta_1)^2}
\right]
= K_2\,\bigl(1 - D\dot\theta_1\bigr),
\end{equation}
with
\begin{equation}
\label{eq:rmt-K2}
K_2 = \sum_s \frac{\omega_s^2}{(1 + D \omega_s \theta_1)^2}.
\end{equation}
Substituting \eqref{eq:rmt-kappa-dot} into \eqref{eq:rmt-theta-dot} and solving the linear equation for $\dot\theta_1$,
\begin{equation}
\label{eq:rmt-upsilon-def}
-D\dot\theta_1
= \frac{\Upsilon}{1 - \Upsilon},
\qquad
\Upsilon := \frac{D K_2 J}{\kappa_1^2}.
\end{equation}
The denominator $1-\Upsilon$ is exactly the Neumann series of the derivative of the fixed-point map; it is positive whenever $\Upsilon < 1$, which is the regime covered by the theorem.

It remains to evaluate $K_2$ in closed form.
With $y = (1-\omega)^s$, $\sum_s \omega^2 (1-\omega)^{2s} f((1-\omega)^s) \to \omega\int_0^1 y f(y)\,dy$, so
\begin{equation}
\label{eq:rmt-K2-integral}
K_2
= \omega\int_0^1 \frac{y}{(1 + \zeta \theta_1 y)^2}\,dy
= \frac{\omega}{(\zeta \theta_1)^2}
\left[\ln(1 + \zeta \theta_1) - \frac{\zeta \theta_1}{1 + \zeta \theta_1}\right].
\end{equation}
Using $e^{\zeta q} = 1 + \zeta \theta_1$ and $\kappa_1^2 = (\zeta q)^2/(\zeta \theta_1)^2$,
\begin{equation}
\label{eq:rmt-K2-collapse}
\frac{D K_2}{\kappa_1^2}
= \frac{\zeta}{(\zeta q)^2}\Bigl[\zeta q - \bigl(1 - e^{-\zeta q}\bigr)\Bigr]
= \frac{1}{q}\left(1 - \frac{1 - e^{-\zeta q}}{\zeta q}\right).
\end{equation}
Combining \eqref{eq:rmt-upsilon-def} and \eqref{eq:rmt-K2-collapse} yields
\begin{equation}
\label{eq:rmt-upsilon}
\Upsilon
= \frac{1}{q}\left(1 - \frac{1 - e^{-\zeta q}}{\zeta q}\right)
\int \left(\frac{\tau}{\tau + \rho_{\mathrm{eff}}}\right)^{2} dH(\tau),
\end{equation}
which is the variance functional of Theorem~\ref{thm:rmt}(i).
The $\zeta \to 0$ limit is consistent: expanding $e^{-\zeta q} = 1 - \zeta q + \tfrac12(\zeta q)^2 + O(\zeta^3)$ gives $\Upsilon = \tfrac12\zeta J + O(\zeta^2)$, so $\sigma^2\Upsilon/(1-\Upsilon) \to \tfrac12\sigma^2\omega\,d_{\mathrm{eff}}(\rho)$, the mean-field variance of Theorem~\ref{thm:window}.

\subsubsection{Pull sensitivities (Theorem~\ref{thm:rmt}(ii))}
\label{app:rmt-sensitivity}

Implicit differentiation of \eqref{eq:rmt-collapse} in $\ln\rho$ gives
\begin{equation}
\label{eq:rmt-sensitivity}
\rho\,\frac{\partial \bar s}{\partial \rho}
= \frac{\mathcal{I}}{1 + \zeta\,\phi(\zeta q)\,\mathcal{I}},
\qquad
\mathcal{I} = \int \tilde s\,(1 - \tilde s)\, dH(\tau),
\end{equation}
with $\tilde s(\tau) = \rho_{\mathrm{eff}}/(\tau + \rho_{\mathrm{eff}})$ and $\phi(u) = (1 - e^{-u})^{-1} - u^{-1} \in (0,1)$.
The anisotropy bound follows from the variance decomposition
\begin{equation}
\label{eq:rmt-anisotropy}
\mathcal{I} = \bar s\,(1 - \bar s) - \operatorname{Var}_H(\tilde s).
\end{equation}
Spectral spread pushes directions into the data-dominated bulk ($\tilde s \approx 0$) or the anchor-pinned tail ($\tilde s \approx 1$), both insensitive, so the closed-loop sensitivity
$S_{\omega\rho} = \mathcal{I}/(q + \zeta\phi(\zeta q)\mathcal{I})$ stays below one on every spectrum.
This is the analytic explanation of the measured coupling margins of Theorem~\ref{thm:closedloop}.

\subsubsection{Conflict bias (Theorem~\ref{thm:rmt}(iii))}
\label{app:rmt-conflict}

Part (i) is the variance component.
The general excess risk adds the bias carried by an anchor offset $\delta = G - w^*$.
Writing the anchored solve as $W = (\bar A + \rho I)^{-1}(\bar B + \rho G)$ and taking expectations over the noise alone, the deterministic part of the deviation is the resolvent action $\rho R(\rho)\delta$ with $R(\rho) = (\bar A + \rho I)^{-1}$.
For \emph{out-of-sample} prediction risk at $x_{\mathrm{new}} \sim (0,\Sigma)$ the conflict contribution is
\begin{equation}
\label{eq:rmt-conflict-quadratic}
\delta^\top \rho^2 R(\rho)\,\Sigma\,R(\rho)\,\delta .
\end{equation}
The in-sample empirical risk replaces $\Sigma$ by $\bar A$ and gives $\mu_i(\rho/(\mu_i + \rho))^2\delta_i^2$ in the classical limit; we state the out-of-sample form, which is the risk the experiments estimate.

In the proportional streaming limit with $\delta$ held fixed, define the direction-weighted conflict measure $H_\delta$ (an atom of mass $\delta_i^2$ at each population direction $\tau_i$, total mass $\|\delta\|^2$).
The deterministic equivalent of \eqref{eq:rmt-conflict-quadratic} is
\begin{equation}
\label{eq:rmt-conflict-de}
\rho_{\mathrm{eff}}^2 \int \frac{\tau}{(\tau + \rho_{\mathrm{eff}})^2}\, dH_\delta(\tau).
\end{equation}
It remains to justify the renormalization by $(1-\Upsilon)^{-1}$.
The perturbation $\varepsilon$ enters \eqref{eq:rmt-conflict-quadratic} only through the scalar $\kappa_1(\varepsilon)$, because the resolvent is a function of $\kappa_1(\varepsilon)\Sigma + \rho I$.
Expanding the resolvent around $\varepsilon = 0$ produces the same scalar Neumann series $1 + \Upsilon + \Upsilon^2 + \cdots$ that produced \eqref{eq:rmt-upsilon-def}, with the same $\Upsilon$ from \eqref{eq:rmt-upsilon}; the conflict quadratic form multiplies that series term by term.
Hence
\begin{equation}
\label{eq:rmt-risk}
R_{\mathrm{pred}}(\rho, \zeta) \;\to\;
\frac{1}{1 - \Upsilon}\left[\sigma^2\,\Upsilon
+ \rho_{\mathrm{eff}}^2 \int \frac{\tau}{(\tau + \rho_{\mathrm{eff}})^2}\, dH_\delta(\tau)\right],
\end{equation}
which is Theorem~\ref{thm:rmt}(iii).
The oracle anchor ($H_\delta \equiv 0$) reduces \eqref{eq:rmt-risk} to Part (i), and the classical limit $\zeta \to 0$ gives $\rho_{\mathrm{eff}} \to \rho$, $\Upsilon \to 0$, recovering Proposition~\ref{prop:conflict}'s per-direction bias.
The derivation is a deterministic-equivalent expansion rather than a finite-$D$ theorem; the rank-deficient case is covered by the corollary below (Appendix~\ref{app:rmt-validation}).

\subsubsection{Active-subspace reduction (rank-deficient spectra)}
\label{app:rmt-active}

The full-rank hypothesis in Theorem~\ref{thm:rmt} is not needed.
Exact null directions carry no prediction risk, and the deterministic equivalent restricts to the range of the population covariance.

\begin{corollary}[Active-subspace deterministic equivalent]
\label{cor:rmt-active}
Let the population spectral measure split as $H = (1-p)\,\delta_0 + p\,H^+$ with $p \in (0,1]$ and $\operatorname{supp}(H^+) \subset [\tau^+_{\min}, \tau^+_{\max}]$, $\tau^+_{\min} > 0$.
Then the deterministic-equivalent system \eqref{eq:rmt-system}--\eqref{eq:rmt-upsilon}, the variance functional $\Upsilon$, and the prediction risk \eqref{eq:rmt-risk}, all evaluated at $(H, \zeta)$, coincide with those evaluated at $(H^+, \zeta^+)$ with the reduced aspect ratio $\zeta^+ = p\,\zeta$.
\end{corollary}

\begin{proof}
The features are $x_s = \Sigma^{1/2} z_s$, so $x_s$ vanishes identically on $\ker\Sigma$ and the accumulator $\bar A$ has an exact zero block there.
In the out-of-sample risk $\E[(x_{\mathrm{new}}^\top(W-w^*))^2] = (W-w^*)^\top\Sigma(W-w^*)$ every error component in $\ker\Sigma$ is multiplied by $\Sigma$, so it contributes nothing.
Algebraically, the atom at zero drops out of the Stieltjes integrals:
\[
\theta_1 = \int \frac{\tau}{\kappa_1\tau+\rho}\,dH = p\,\theta_1^+,
\qquad \theta_1^+ := \int \frac{\tau}{\kappa_1\tau+\rho}\,dH^+,
\]
while the weight sum \eqref{eq:rmt-kappa} depends on $\theta_1$ only through the product
$\zeta\,\theta_1 = (D\omega)(p\,\theta_1^+) = (pD)\,\omega\,\theta_1^+ = \zeta^+ \theta_1^+$.
Hence \eqref{eq:rmt-system} is the $H^+$ system at aspect ratio $\zeta^+$, so $\kappa_1$ and therefore $\rho_{\mathrm{eff}} = \rho/\kappa_1$ are unchanged.
Retentions scale as $q = \kappa_1\theta_1 = p\,q^+$ with the active retention $q^+ := \kappa_1\theta_1^+ = \int \tau/(\tau+\rho_{\mathrm{eff}})\,dH^+$, which is the fixed point \eqref{eq:rmt-collapse} with $\rho_{\mathrm{eff}} = \rho (e^{\zeta^+ q^+}-1)/(\zeta^+ q^+)$.
For the variance functional, $J = \int (\tau/(\tau+\rho_{\mathrm{eff}}))^2\,dH = p\,J^+$ and $D\,K_2/\kappa_1^2$ is a function of $\zeta\theta_1 = \zeta^+\theta_1^+$ alone, so
$\Upsilon = (D K_2/\kappa_1^2)\,J = (D^+ K_2/\kappa_1^2)\,J^+ = \Upsilon^+$.
The conflict integral in \eqref{eq:rmt-risk} carries the weight $\tau/(\tau+\rho_{\mathrm{eff}})^2$, which vanishes on the null atom, so it is unchanged as well.
\end{proof}

The corollary says a rank-deficient population is not a boundary of the theory: it is the same deterministic equivalent on the active subspace, of dimension $D^+ = pD$ and aspect ratio $\zeta^+ = p\zeta$, with the hypothesis reduced to $\tau_{\min}^+ > 0$ on the nonzero part.
The direction-weighted conflict measure $H_\delta$ is likewise supported on $\operatorname{range}(\Sigma)$ in the prediction risk; the unweighted trace functional that ignores the $\tau$ weighting is the object that fails, not the theorem.

\subsubsection{The coupling corollary is local}
\label{app:rmt-coupling}

At an interior fixed point the product of cross-sensitivities is invariant under monotone coordinate changes, since $\partial \psi_1(\omega^*)/\partial \psi_2(\rho^*)\cdot \partial \psi_2(\rho^*)/\partial \psi_1(\omega^*) = |(\partial \omega^*/\partial \rho)(\partial \rho^*/\partial \omega)|$; we therefore evaluate in $(\ln \omega, \ln \rho)$ without loss of generality.
(i) The window target $\omega^* \propto (4\nu^2_\Sigma / c')^{1/3}$ with $c' = \sigma^2 d_{\mathrm{eff}}(\rho) \propto q^2$ gives $|\partial \ln \omega^*/\partial \ln \rho| = \tfrac{2}{3}\,\mathcal I / [q(1 + \zeta\phi(\zeta q)\mathcal I)] \le \tfrac{2}{3}\,\bar s/q$, using $\mathcal I \le \bar s(1-\bar s) = q\bar s$.
(ii) The strength target $\rho^*(\omega) = \bar\mu\widehat V(\omega)/\widehat C$ is analyzed by which branch of $\widehat C$ is active:
(a) when $C^\Sigma_{\mathrm{floor}}$ is active, $\widehat C$ is independent of $\omega$, and $|\partial \ln \rho^*/\partial \ln \omega| \le 2/(2-\omega)$;
(b) when $C^\Sigma_{\mathrm{resid}} = E_t$ is active, the drift-active floor $\rho_{\min}^{\mathrm{drift}}$ pins $\rho$, so $L_{\rho\omega} = 0$ and the system is triangular (Theorem~\ref{thm:closedloop}(ii));
(c) when $C^\Sigma_{\mathrm{ref}}$ is active at stationarity ($\nu = 0$), $\widehat w$ is unbiased for $w^*$ in the mean-field limit, so the reference branch is independent of $\omega$ to leading order in $1/n$ and $L_{\rho\omega} = 0$; away from stationarity it inherits the same $(\omega/(2-\omega))$ scaling as (a), with the constant absorbed into $\bar s$.
Branches (b) and (c) carry explicit finite-$n$ corrections. In branch (b) the drift-active floor fixes $\rho$ exactly, so $L_{\rho\omega} = 0$; the residual dependence is only the branch-selection margin, bounded by the $O(1/n)$ fluctuation of the residual excess $E_t$.
In branch (c), under the quasi-static estimator block of Theorem~\ref{thm:closedloop}, the computed quadratic form $R = (G_{\mathrm{ref}}-\widehat w)^\top \bar A_t (G_{\mathrm{ref}}-\widehat w)$ is $\omega$-independent, so the runtime estimate $\widehat C_{\mathrm{ref}}(\omega) = R - \widehat V(\omega)$ has the single explicit $\omega$-dependence of the debiasing term $\widehat V(\omega) = \widehat\sigma^2 d_{\mathrm{act}}\,\omega/(2-\omega)$; differentiating $\rho^* = \bar\mu\widehat V/\widehat C_{\mathrm{ref}}$,
\[
\bigl|L_{\rho\omega}^{(c)}\bigr|
= \frac{2}{2-\omega}\Bigl(1 + \frac{\widehat V}{|\widehat C_{\mathrm{ref}}|}\Bigr),
\]
the leading coefficient plus the measurable debiasing ratio $\widehat V/|\widehat C_{\mathrm{ref}}| = \rho^*/\bar\mu$ at the fixed point, in place of an unbound $O(1/n)$.
The measured coupling product is below $6\times10^{-3}$ across all $20$ diagnostic runs, with one quantity pinned in three of four configurations (\texttt{results/coupling\_diag.csv}).
The coupling product in log-coordinates is therefore at most
\begin{equation}
\label{eq:rmt-coupling}
L_{\ln \omega, \ln\rho}\,L_{\ln\rho, \ln \omega}
\;\le\; \frac{2}{3q}\,\mathcal I \cdot \frac{2}{2-\omega}
\;=\; \frac{4\,\mathcal I}{3q(2-\omega)}
\;\le\; \frac{4}{3(2-\omega)}\,\bar s
\;<\; \bar s
\;<\; 1
\end{equation}
for all $\bar s \in (0,1)$ and $\omega \in (0, 2/3)$, with strict improvement from spectral anisotropy ($\operatorname{Var}_H(\tilde s) > 0$ reduces $\mathcal I$ below $q\bar s$).
This is a \emph{local} contraction statement at an interior fixed point with the clip inactive: it bounds the Jacobian of the estimator map, not the iterates after a branch switch or while a cap is active.
We therefore claim local contraction, not global contraction along every trajectory; the global statement in Theorem~\ref{thm:closedloop} retains the measured-product hypothesis and the gate.

\subsubsection{Numerical validation}
\label{app:rmt-validation}

The closed forms above were checked against direct simulation (\texttt{scripts/diagnostics.py} \texttt{--diag rmt}), separately from the proof.
\emph{Fixed points.}
At $45$ operating points (three spectra $\times$ five nominal windows $1/(1-\lambda) \in \{100,\dots,2000\}$ $\times$ three anchor strengths $\rho \in \{10^{-3}, 10^{-2}, 10^{-1}\}$), \eqref{eq:rmt-collapse} matches the Monte Carlo stationary risk to within $4.4\%$ (median $1.4\%$) across all three spectra, and within $2.3\%$ on the rank-deficient spectrum.
The rank-deficient grid is recorded in \texttt{results/rmt\_rankdef.csv}, which also verifies the active-subspace invariance of Corollary~\ref{cor:rmt-active}: the full-spectrum and reduced-system functionals agree to machine precision ($|\Upsilon-\Upsilon^+| = 0$ at every point).
\emph{Sensitivities.}
Equation~\eqref{eq:rmt-sensitivity} agrees with numerical differentiation of the fixed point to $10^{-6}$ relative error across all nine spectrum--operating-point combinations.
\emph{Conflict increment.}
Paired Monte Carlo (same seeds, conflict-anchor total minus oracle total) matches \eqref{eq:rmt-risk} within $0.6$--$6.7\%$ (median $1.1\%$) for a random conflict direction, and within $1$--$9.5\%$ (median $2.2\%$) for a top-eigendirection conflict, across all three spectra including the rank-deficient one; the direction weighting by $\tau/(\tau+\rho_{\mathrm{eff}})^2$ is what keeps the conflict measure correct there.
The random-orientation rank-deficient case is covered by Corollary~\ref{cor:rmt-active}: restricting the deterministic equivalent to $\operatorname{range}(\Sigma)$ removes the null-space mismatch, after which the paired Monte Carlo matches \eqref{eq:rmt-risk} within $6.7\%$ at every point.
\emph{Operating points.}
At the method's own converged control pair the deterministic equivalent matches the arm's Monte Carlo risk to within $9.8\%$ (isotropic), $12.7\%$ (decay3), and $1.9\%$ (rank-deficient).
\emph{Finite-$D$ concentration.}
Scoring the risk with the population covariance---the out-of-sample quantity Theorem~\ref{thm:rmt} is stated for---the relative discrepancy $|R_{\mathrm{MC}}-R_{\mathrm{DE}}|/R_{\mathrm{DE}}$ at fixed $\zeta$ is within the Monte Carlo standard error at every $D \in \{64,128,256,512\}$ over $\zeta \in \{1,2\}$ and $\rho \in \{10^{-2},10^{-1}\}$: for $D \ge 128$ the maximum relative bias is $1.1\%$ and the median $|\text{bias}|/\text{s.e.}$ is $0.66$, with no systematic growth in $D$, consistent with an $O(1/D)$ deterministic-equivalent error (\texttt{py scripts/finite\_d\_conc.py}; \texttt{results/finite\_d\_conc.csv}).
The fixed-point grid above scores the risk with the in-sample empirical covariance, which adds an $O(D/T)$ in-sample term ($\approx 1.3\%$ at $D=128$, $T=10^4$) on top of this finite-$D$ error.
The CSV files \texttt{results/rmt\_grid.csv}, \texttt{results/rmt\_verify.csv}, \texttt{results/rmt\_conflict.csv}, and \texttt{results/rmt\_rankdef.csv} record the grids; all rows, full-rank and rank-deficient, support the theorem.

The convergence step is cited \citep{silverstein1995,rubio2011,hachem2007}; the nearest prior analysis is the large-system exponentially windowed RLS study of \citet{xiao2005}, without anchoring, per-direction pulls, refresh, or self-calibrating estimators, and the unweighted-batch ridge asymptotics \citep{hastie2022,pafka2004} do not carry the finite-window renormalization.

\subsection{Closed-loop stability of the two estimators (Theorem~\ref{thm:closedloop})}

\emph{Part (i): the coupled recursions.}
The recursion is the quasi-static two-time-scale model: the estimator block runs every $k$ steps and is smoothed with horizons $1/\eta_C, 1/\eta_\nu$, both slow relative to the accumulator mixing time $1/\omega$ and the fast EMA horizon $h_f$; finite-$k$ lag is an additive perturbation absorbed into $\xi_\omega, \xi_\rho$ \citep{borkar2008,benveniste1990,kushner2003}.
The window estimator maintains an exponential average of the debiased drift-energy estimate, $\widehat\nu^2_{m+1} = (1-\eta_\nu)\widehat\nu^2_m + \eta_\nu\,\widehat\nu^2_\Sigma(\rho_m)$ with gain $\eta_\nu$ (the estimator of Theorem~\ref{thm:window}); at a frozen target the cubed readout $u = \omega^3 = 4\widehat\nu^2/c'$ is affine in the averaged quantity, and the frozen-target contraction is formalized (\texttt{debiased\_dial\_contraction}).
Since $u$ is affine in $\widehat\nu^2$ at frozen $c'$, and $c' = (1-\bar s)^2\sigma^2D$ is Lipschitz in $\rho$ through $\bar s$, the error $e_\omega(m) = |u_m - (\omega^*(\rho_m))^3|$ satisfies
\[
e_\omega(m{+}1) \;\le\; (1-\eta_\nu)\,e_\omega(m) + \eta_\nu L_{\omega\rho}\, e_\rho(m) + \xi_\omega(m),
\]
where $L_{\omega\rho}$ absorbs the readout and $c'$ sensitivities and $\xi_\omega$ collects the drift-energy estimation noise, bounded by the same sub-exponential concentration as Proposition~\ref{prop:falsealarm}.
The clip onto $[\omega_{\min}, \omega_{\max}]$ is a projection onto a convex set, which preserves the one-sided error contraction; a pinned estimator is the triangular case of part~(ii).
The strength estimator tracks the plug-in target $\rho^*(\omega) = \bar\mu\widehat V/\widehat C$ of Proposition~\ref{prop:conflict} through clipping and shrinkage; $\widehat V$ and $\widehat C$ are Lipschitz in the window coordinate on any compact parameter box that excludes eigenvalue crossings of the $10^{-8}$ ridge: $d_{\mathrm{act}}$ is integer-valued and jumps only when an eigenvalue crosses that ridge, so we either assume a spectral gap (the grid's spectra satisfy it with margin) or replace $d_{\mathrm{act}}$ by the smooth effective dimension $\mathrm{tr}(\bar A(\bar A+\rho_{\mathrm{ridge}}I)^{-1})$; on that set, through $\bar\mu(\omega)$, $d_{\mathrm{act}}(\omega)$, and the prediction-unit reference conflict $(G_{\mathrm{ref}}-\widehat w)^\top\bar A(G_{\mathrm{ref}}-\widehat w)$, the maps are Lipschitz, giving
\[
e_\rho(m{+}1) \;\le\; \eta_C L_{\rho\omega}\, e_\omega(m) + (1-\eta_C)\,e_\rho(m) + \xi_\rho(m).
\]
Setting $a = \eta_\nu L_{\omega\rho}$, $b = \eta_C L_{\rho\omega}$, the two recursions form the coupled system of \texttt{Coupled.lean}.

\emph{The weighted-norm contraction.}
With the midpoint weight $r := \tfrac12(a/\eta_\nu + \eta_C/b)$---the interval $(a/\eta_\nu,\ \eta_C/b)$ is nonempty exactly when $ab < \eta_\nu\eta_C$, i.e.\ $L_{\omega\rho}L_{\rho\omega} < 1$---the joint rate is
\[
\gamma_{\mathrm{cl}} \;=\; \max\bigl\{(1-\eta_\nu) + a/r,\ \ (1-\eta_C) + r\,b\bigr\} \;<\; 1,
\]
and the error pair contracts in the weighted max-norm $\max\{e_\omega,\, r\,e_\rho\}$: the one-step bound, geometric decay, and tracking radius $\xi/(1-\gamma)$ are formalized as \texttt{joint\_step}, \texttt{joint\_geometric}, and \texttt{joint\_tracking}, with the rate lemma \texttt{joint\_rate\_lt\_one} and the weight lemmas \texttt{joint\_bounds}, \texttt{joint\_weight\_lo}, \texttt{joint\_weight\_hi}.

\emph{Part (ii): the drift-mode pin.}
When the detector's drift-active floor pins $\rho$, $e_\rho \equiv 0$ and the window error contracts at the free estimator's own rate $1-\beta$ (\texttt{joint\_geometric\_pinned}); the joint-rate specialization is \texttt{joint\_geometric\_drift\_mode}.
The complementary one-way case (the strength target insensitive to the window) is the triangular specialization \texttt{joint\_geometric\_triangular}.

\emph{Part (iii): the fluctuation bound.}
The contracted recursion of part (i) is the scalar Lyapunov inequality whose stationary second moment is bounded in Proposition~\ref{prop:fluct}; the matrix second-moment route is deliberately not claimed, because entrywise domination of nonnegative error vectors does not lift to the Loewner order (a nonnegative triangle-inequality remainder also cannot be a zero-mean martingale difference), so the scalar weighted-norm route is the one the contraction actually supports.

\begin{proposition}[Invariant measure in the co-moving frame]
\label{prop:ergodic}
Assume $(x_t,\varepsilon_t)$ is i.i.d.\ (or stationary ergodic) with $\E\|x_t\|^2 \le \bar\Sigma$, $\varepsilon_t$ sub-Gaussian, and constant drift $\nu$.
Work in the co-moving frame $\tilde B_t := B_t - A_t w^*_t$ and $\tilde G_t := G_t - w^*_t$, and let $\Phi_m$ be the $k$-step skeleton at $t = mk$.
\begin{enumerate}
\item[(i)] $\Phi_m$ is a Feller Markov chain on the product of the compact parameter box with the trace-bounded (tight) accumulator set, and admits an invariant probability measure $\pi$.
\item[(ii)] \sloppy If the mean-field estimator map of Theorem~\ref{thm:closedloop} is contracting ($\gamma_{\mathrm{cl}}<1$) and the full state-to-estimator Jacobian satisfies the small-gain condition $\varrho(J_{\mathrm{state}\to\mathrm{estimator}})\,\varrho(J_{\mathrm{estimator}\to\mathrm{state}}) < 1-\gamma_{\mathrm{cl}}$, then $\pi$ is unique and $\Phi_m$ is geometrically ergodic in $W_1$.
\end{enumerate}
No pathwise synchronous-coupling claim is made: when the two copies' estimators differ, the accumulator difference carries the cross term $-(\omega^{(1)}-\omega^{(2)})A^{(2)}$, which the small-gain condition controls.
\end{proposition}
\begin{proof}
(i) The co-moving frame removes the linear growth: $A_t$ is tight by the Foster--Lyapunov drift $\E[A_{t+1}\mid\mathcal F_t]\preceq(1-\omega_{\min})A_t+\bar\Sigma$, the centered statistics $\tilde B_t,\tilde G_t$ are bounded through the pull interpolation, and the estimators live in a compact box by construction.
The one-step kernel is $k$-periodic because the anchor refresh and estimator updates occur every $k$ steps, so Krylov--Bogoliubov is applied to the time-homogeneous $k$-step skeleton $\Phi_m$.
The hard gate $\mathbf 1\{E_t>0\}$ does not break the Feller property: when $\varepsilon_t$ has a continuous density (Gaussian or bounded with a density), the boundary event $\{E_t=0\}$ has conditional probability zero, so the transition expectation of any bounded continuous function is continuous in the state; no artificial smoothing of the gate is needed.
Tightness plus this Feller property yields an invariant probability measure \citep{meyn2009}.
(ii) Linearizing the joint error dynamics around the invariant set gives the feedback system state $\to$ estimator $\to$ state. The small-gain condition makes its spectral radius strictly below one in a weighted norm, and combined with the Foster--Lyapunov drift and the support condition this gives geometric ergodicity and uniqueness \citep[Ch.~18]{meyn2009}.
Without the small-gain condition, uniqueness is not claimed for the full process; the idealized mean-field estimator process of Theorem~\ref{thm:closedloop} is unique under $\gamma_{\mathrm{cl}}<1$.
\end{proof}

\begin{proposition}[Fluctuation bound]
\label{prop:fluct}
Let the coupled estimator deviations obey the weighted max-norm recursion of part (i), $V(m{+}1) \le \gamma_{\mathrm{cl}}\,V(m) + \|\xi(m)\|_r$, where $\xi(m)$ is any perturbation sequence with uniformly bounded second moment $\sup_m \E\|\xi(m)\|_r^2 \le q_*$.
Then the stationary second moment satisfies $\limsup_m \E[V(m)^2] \le q_*\,/\,(1 - \gamma_{\mathrm{cl}})^2$; in particular each estimator fluctuates within RMS distance $\sqrt{q_*}\,/\,(1 - \gamma_{\mathrm{cl}})$ of its target.
No zero-mean, martingale, or independence assumption is used.
\end{proposition}
\begin{proof}
Young's inequality $(a + b)^2 \le (1 + \epsilon)a^2 + (1 + 1/\epsilon)\,b^2$ with $\epsilon = 1/\gamma_{\mathrm{cl}} - 1$ makes $(1 + \epsilon)\,\gamma_{\mathrm{cl}}^2 = \gamma_{\mathrm{cl}}$ and $(1 + 1/\epsilon) = 1/(1 - \gamma_{\mathrm{cl}})$, so
\[
  \E[V(m{+}1)^2] \;\le\; \gamma_{\mathrm{cl}}\,\E[V(m)^2] \;+\; \E\|\xi(m)\|_r^2\,/\,(1 - \gamma_{\mathrm{cl}}) \;\le\; \gamma_{\mathrm{cl}}\,\E[V(m)^2] + q_*/(1 - \gamma_{\mathrm{cl}}).
\]
Iterating from a bounded initial second moment and summing the geometric series $\sum_j \gamma_{\mathrm{cl}}^j = 1/(1 - \gamma_{\mathrm{cl}})$ gives $\E[V(m)^2] \le \gamma_{\mathrm{cl}}^m\,\E[V(0)^2] + q_*/(1 - \gamma_{\mathrm{cl}})^2$.
\end{proof}
No zero-mean or martingale assumption is used: Young's inequality bounds the cross term $2ab$ pathwise. If $\xi$ were a signed martingale difference and one worked with the signed recursion before taking norms, the cross term would vanish and the denominator would improve from $(1-\gamma_{\mathrm{cl}})^2$ to $1-\gamma_{\mathrm{cl}}^2$; the norm route is deliberately more general and covers the nonnegative triangle-inequality remainder as well.

This closes the loop: the two single-estimator analyses are the frozen-coupling cases, the product condition is the coupling budget, and the gate enforces it exactly where the targets move fastest.
The deterministic coupled stationary law is the unique fixed point of the composition $\Phi = \rho^*\circ\omega^*$: substituting the window target $\omega^*(\rho) = (4\nu^2_\Sigma/\sigma^2 d_{\mathrm{eff}}(\rho))^{1/3}$ into the strength target $\rho^*(\omega) = \bar\mu\widehat V(\omega)/\widehat C(\omega)$ gives the scalar equation
\[
\rho \;=\; \Phi(\rho), \qquad
\Phi(\rho) \;=\; \frac{\bar\mu\,\widehat V(\omega^*(\rho))}{\widehat C(\omega^*(\rho))},
\]
a contraction with modulus the coupling product $L_{\omega\rho}L_{\rho\omega} < 1$; hence $\Phi$ has a unique fixed point, reached from any start by the iteration $\rho_{m+1} = \Phi(\rho_m)$, and the law admits the convergent telescoping series $\rho^* = \sum_{m\ge0}(\rho_{m+1}-\rho_m)$ with geometric remainder $O((L_{\omega\rho}L_{\rho\omega})^m)$.
No single elementary closed form is available: $\Phi$ composes a cube root of the rational $d_{\mathrm{eff}}(\rho)$ with the rational strength map, so the fixed-point equation is transcendental and the convergent series is the closed-form handle.
The stochastic stationary law's mean is this fixed point and its second moment is bounded by $q_*/(1-\gamma_{\mathrm{cl}})^2$ (Proposition~\ref{prop:fluct}); the moment recursion is the exact form available, and no closed-form density is claimed because the innovation $\xi$ is state dependent.

\subsection{Spectral Invariance of the Scalar Method}
\label{app:invariance}

\begin{proposition}[Scalar-method rotation equivariance]
\label{prop:invariance}
For any orthogonal $Q$, the scalar-$\rho$ anchored solve is equivariant under rotating the stream, $x'_t = Qx_t$ with anchor $G' = QG$: the rotated solve is $W' = QW$, every prediction is invariant ($x'^\top_t W' = x^\top_t W$), and the spectral refresh commutes with the rotation.
\end{proposition}
\begin{proof}
The rotated accumulators are $A' = QAQ^\top$ and $B' = QB$, and $QAQ^\top + n\rho I = Q(A + n\rho I)Q^\top$, so
$W' = (QAQ^\top + n\rho I)^{-1}(QB + n\rho\,QG) = Q\,(A + n\rho I)^{-1}(B + n\rho G) = QW$.
Predictions are inner products of rotated vectors against rotated solutions.
The eigenvalues of $\bar A' = Q\bar A Q^\top$ coincide with those of $\bar A$, so the pulls $s_i = \rho/(\mu_i + \rho)$ are unchanged; with the rotated eigenbasis $U' = QU$ the refresh step maps to the rotated refresh, so the equilibrium of Theorem~\ref{thm:cadence} is unchanged as well.
The trace and determinant invariance of the conjugated accumulator are formalized as \texttt{conj\_trace} and \texttt{conj\_det} in \texttt{Core.lean}.
\end{proof}
Two consequences: the scalar method's per-step solve never needs an eigenbasis (it uses only $A$, $B$, $G$), so a covariance rotation mid-stream degrades it through the conditioning of $A + n\rho I$, which Proposition~\ref{prop:windup} bounds, and through the frozen anchor's misalignment until the next refresh; and the per-direction method's refresh aligns to the current dominant directions only at each cycle, so between cycles it applies its pulls in a stale basis, at the Davis--Kahan-bounded cost above. The scalar method wins all nine covariance-rotation configurations while the per-direction method is the fragile one (Sections~\ref{sec:experiments} and~\ref{sec:limitations}).

\section{Supporting Theory}
\label{app:support}

This section collects the supporting results behind the anchored solve and the effective-memory-horizon analysis, with short proofs and references to the standard literature.

\subsection{The Anchored Solve as Exact MAP}
\label{app:support-map}

For a forgetting factor $\alpha \in (0,1]$ define the weighted statistics
\[
  A_t = \sum_{i=1}^t \alpha^{t-i}\,\phi_i\phi_i^\top, \qquad
  B_t = \sum_{i=1}^t \alpha^{t-i}\,\phi_i y_i^\top, \qquad
  M_t = \sum_{i=1}^t \alpha^{t-i},
\]
and the anchored objective with prior mean $G$, prior precision $\tau_0 > 0$, and ridge $\lambda \ge 0$,
\[
  J_t(W) \;=\; \sum_{i=1}^t \alpha^{t-i}\,\|\phi_i^\top W - y_i\|_F^2
  \;+\; (\tau_0 + \lambda)\,\|W - G\|_F^2 .
\]
The anchor penalty is undiscounted because $G$ is a fixed prior, not streaming data.

\begin{proposition}[Anchored ridge as exact MAP]
\label{prop:app-map}
Assume $y_i = \phi_i^\top W^\star + \varepsilon_i$ with $\varepsilon_i \sim \mathcal N(0, \sigma^2 I)$ and prior $W \sim \mathcal N(G, \tau^{-1} I)$ with $\tau = (\tau_0 + \lambda)/\sigma^2$.
Then the posterior mean of $W$ given $(\phi_i, y_i)_{i \le t}$ equals
\[
  W_t^\star \;=\; \bigl(A_t + (\tau_0 + \lambda) I\bigr)^{-1}\bigl(B_t + (\tau_0 + \lambda) G\bigr).
\]
\end{proposition}
\begin{proof}
Completing the square in $J_t$: the normal equations are $(A_t + (\tau_0 + \lambda)I)\,W = B_t + (\tau_0 + \lambda)G$ with positive-definite Hessian, and the Gaussian posterior mean solves the same system.
\end{proof}

The identity is the standard ridge-as-MAP correspondence \citep{hoerl1970,rasmussen2006}, with regularization toward a nonzero prior mean as in L2-SP \citep{li2018l2sp}.
Two algebraic facts are used in the text: the solve is ordinary ridge on the augmented design $[\Phi;\ \sqrt{\tau_0 + \lambda}\,I]$ with pseudo-targets $[Y;\ \sqrt{\tau_0 + \lambda}\,G]$, so standard conditioning statements transfer mechanically and the $\tau_0 \to \infty$ limit is exactly the frozen warm solution (the zero-shot endpoint of the WiSE-FT interpolation \citep{wortsman2022wiseft}); and the count-scaled form $W = (A_t + M_t\rho I)^{-1}(B_t + M_t \rho G)$ used by SPAR matches \hyperref[prop:app-map]{the proposition} with the anchor weight proportional to the effective count, whereas an anchor weight that does not scale with the accumulated count dilutes as the stream grows, biasing strength sweeps.

\subsection{Effective-memory-horizon Stability and the Rank Obstruction}
\label{app:support-emh}

\begin{proposition}[Unanchored stability needs excitation over the window]
\label{prop:app-emh}
Let $\alpha = 1 - 1/m$ with integer $m \ge 2$, and suppose every window $\mathcal W_t = \{t - m + 1, \dots, t\}$ satisfies the excitation condition
\[
  \lambda_{\min}\Bigl(\sum_{i \in \mathcal W_t} \phi_i \phi_i^\top\Bigr) \;\ge\; s^2 \;>\; 0 .
\]
If $\|\phi_i\| \le L$ for all $i$, then
\[
  \lambda_{\min}(A_t) \;\ge\; e^{-1} s^2,
  \qquad
  \|A_t^{-1}\| \;\le\; e/s^2,
  \qquad
  \kappa_2(A_t) \;\le\; e\,m\,L^2/s^2 .
\]
\end{proposition}
\begin{proof}
For samples in the window, $\alpha^{t-i} \ge \alpha^{m-1} = (1 - 1/m)^{m-1} = \bigl(1 + \tfrac{1}{m-1}\bigr)^{-(m-1)} \ge e^{-1}$, since $(1 + \tfrac1n)^n \le e$ for all $n \ge 1$ and taking reciprocals gives the required lower bound.
Hence $A_t \succeq e^{-1} \sum_{i \in \mathcal W_t} \phi_i\phi_i^\top \succeq e^{-1} s^2 I$.
For the upper bound, $\lambda_{\max}(A_t) \le L^2 \sum_{k \ge 0} \alpha^k = mL^2$ (with $A_t$ truncated at $t$), which gives the condition-number bound.
\end{proof}

\begin{corollary}[Rank obstruction]
\label{cor:app-rank}
The window sum $\sum_{i \in \mathcal W_t} \phi_i\phi_i^\top$ has rank at most $\min(m, D)$.
Hence if $m < D$ the excitation condition can hold only on a subspace of dimension $m$, and in the generic case the unanchored recursion is ill-posed when $m \ll D$.
This is the covariance windup of aggressive forgetting, whose measured boundary is reported in Section~\ref{sec:experiments}.
\end{corollary}

\begin{corollary}[Anchoring removes the excitation requirement]
\label{cor:app-anchor-excites}
For any $\tau_0 + \lambda > 0$,
$A_t + (\tau_0 + \lambda)I \succeq (\tau_0 + \lambda)I$ regardless of excitation, so
$\|(A_t + (\tau_0 + \lambda)I)^{-1}\| \le 1/(\tau_0 + \lambda)$ and
$\kappa_2 \le 1 + mL^2/(\tau_0 + \lambda)$: the anchor replaces the persistence-of-excitation assumption by a conditioning budget.
\end{corollary}

The practical design rule is $m = 1/\omega \gtrsim D$ (the effective window at least the dimension), or an anchor strength of at least $10^{-2}$ when $m < D$.
Proposition~\ref{prop:app-emh} and its corollaries are the stability statement behind Proposition~\ref{prop:windup} of the main text: the anchored solve never relies on excitation, and the forgetting method that violates the rule is the one that diverges.

\section{The Lean 4 Artifact}
\label{app:lean}


The algebraic core and the no-regret chain of Propositions~\ref{prop:conflict}--\ref{prop:ocowrap} are formalized in Lean 4 with Mathlib (\texttt{theory/lean/SparTheorems/} in the artifact).
This section maps the paper's results to the artifact, states the formalization boundary, records the axiom dependencies, and gives the build recipe.

\subsection{Result-to-declaration Map}
Table~\ref{tab:lean-map} summarizes the correspondence between the paper's results and the Lean 4 declarations.

\begin{table}[!htbp]
\centering
\scriptsize
\setlength{\tabcolsep}{3pt}
\caption{Map from the paper's results to the Lean artifact.
The differentiation step of Theorem~\ref{thm:window}'s convexity and the probabilistic statements are pen-and-paper; every algebraic identity with a named Lean declaration above is formalized; the remaining rows are pen-and-paper and appear in the list below.}
\label{tab:lean-map}
\resizebox{\textwidth}{!}{%
\begin{tabular}{l p{1.2in} p{1.9in} l}
\toprule
paper result & content & Lean declaration(s) & file \\
\midrule
Equation~\eqref{eq:pull} & pull identity ($\gamma = 1$ and general) & \texttt{pull\_identity}, \texttt{pull\_identity\_general} & Core.lean \\
Equation~\eqref{eq:pull} (eigenbasis) & diagonal decoupling; coordinatewise anchored solve & \texttt{eigenbasis\_solve\_coord}, \texttt{eigenbasis\_pull\_identity} & Core.lean \\
Proposition~\ref{prop:conflict} & risk interpolation; dominance; calibration identity; optimal pull; multiplicative inversion at the optimum & \texttt{convex\_interp\_sq}, \texttt{convex\_interp\_sq\_le}, \texttt{dominance}, \texttt{dominance\_strict}, \texttt{quadratic\_calibration}, \texttt{quadratic\_calibration\_delta}, \texttt{optimal\_pull\_min}, \texttt{optimal\_pull\_attained}, \texttt{pull\_inverse} & Core.lean \\
Proposition~\ref{prop:minimax} & minimax value (core inequality); safety excess identity; multiplicative safety threshold $s(v+c)\le 2v$ & \texttt{minimax\_lower}, \texttt{quadratic\_calibration}, \texttt{safety\_threshold} & Core.lean \\
Lemma~\ref{lem:zerolag} & lag-invariance algebra & \texttt{lag\_invariance} & Core.lean \\
Theorem~\ref{thm:cadence} & refresh rate identity; finite-cadence fixed point; geometric closed form; contraction; correlated variance identities and the lower endpoint $\Psi\ge(1-s)^2$ & \texttt{refresh\_rate\_identity}, \texttt{refresh\_fixed\_point}, \texttt{refresh\_closed\_form}, \texttt{refresh\_contraction\_rate}, \texttt{refresh\_anchor\_variance}, \texttt{refresh\_variance\_form}, \texttt{refresh\_variance\_lt}, \texttt{refresh\_variance\_lb} & Core.lean \\
Theorem~\ref{thm:cadence} (moving frame) & moving-frame reduction to the affine recursion; midpoint cycle assembly; Gauss sum & \texttt{moving\_frame\_reduction}, \texttt{cycle\_midpoint\_average}, \texttt{cycle\_lag\_assembly}, \texttt{sum\_range\_half} & Core.lean \\
Theorem~\ref{thm:window} & leading-order window optimum via exact factorization; second-derivative positivity; estimator fixed-point dichotomy & \texttt{window\_optimum}, \texttt{window\_second\_deriv\_pos}, \texttt{naive\_dial\_no\_fixed\_point}, \texttt{debiased\_dial\_contraction} & Ogd.lean, Core.lean \\
Theorem~\ref{thm:window} (variance coefficient) & exact effective dimension; plug-in gap $D(1-\bar s)^2 \le d_{\mathrm{eff}}$ & \texttt{deff\_sub\_plugin}, \texttt{deff\_sub\_plugin\_nonneg} & Core.lean \\
Proposition~\ref{prop:invariance} & rotation equivariance; refresh commutation & \texttt{conj\_trace}, \texttt{conj\_det} (spectral invariance under orthogonal conjugation) & Core.lean \\
Proposition~\ref{prop:falsealarm} & concentration step & --- (pen-and-paper, \citep{freedman1975}) & --- \\
 & subgradient identity & \texttt{sq\_subgradient} & Ogd.lean \\
 & box-projection contraction (1-dim) & \texttt{projBox\_dist\_sq\_le} & Ogd.lean \\
Proposition~\ref{prop:ocowrap}(i) & Zinkevich step (adaptive losses) & \texttt{ogd\_regret\_box}, \texttt{ogd\_regret\_segment} & Ogd.lean \\
 & $D$-dimensional contraction & \texttt{projBoxD\_dist\_sq\_le}, \texttt{sum\_expand\_sq}, \texttt{ogd\_step\_D}, \texttt{ogd\_regret\_box\_D} & BoxD.lean \\
Proposition~\ref{prop:ocowrap}(ii) & quadratic exponential bound & \texttt{exp\_neg\_le\_one\_sub\_add\_sq} & Hedge.lean \\
 & potential step, ratio, and product bounds & \texttt{Zfun\_step}, \texttt{Zfun\_step\_bound}, \texttt{Zfun\_prod\_bound} & Hedge.lean \\
 & Hedge regret & \texttt{hedge\_regret} & Hedge.lean \\
Proposition~\ref{prop:windup} & eigenvalue interval bound; pull weights & \texttt{windup\_ratio\_le}, \texttt{pull\_weight\_bounds} & Core.lean \\
Theorem~\ref{thm:closedloop} & joint $2\times2$ contraction; geometric decay; tracking radius; drift-mode, one-way, and pinned-estimator (free-rate) specializations & \texttt{joint\_step}, \texttt{joint\_geometric}, \texttt{joint\_tracking}, \texttt{joint\_bounds}, \texttt{joint\_weight\_lo}, \texttt{joint\_weight\_hi}, \texttt{joint\_rate\_lt\_one}, \texttt{joint\_geometric\_drift\_mode}, \texttt{joint\_geometric\_triangular}, \texttt{joint\_geometric\_pinned} & Coupled.lean \\
Propositions~\ref{prop:ergodic},~\ref{prop:fluct} & invariant measure; second-moment (Lyapunov) bound & --- (pen-and-paper, \citealp{meyn2009}) & --- \\
\midrule
Corollary~\ref{cor:mixed} & mixed-drift master law & --- (pen-and-paper, Appendix~\ref{app:proofs}) & --- \\
Theorem~\ref{thm:rmt} & spectral retention and streaming risk; collapse/renormalization identities & \texttt{de\_log\_identity}, \texttt{de\_collapse} (convergence pen-and-paper) & Core.lean \\
Props.~\ref{prop:app-map},~\ref{prop:app-emh}; Cors.~\ref{cor:app-rank},~\ref{cor:app-anchor-excites} & exact MAP; excitation/rank; anchoring removes the excitation requirement & --- (pen-and-paper, Appendix~\ref{app:support}) & --- \\
\bottomrule
\end{tabular}}
\end{table}

\subsection{Coverage and Formalization Boundary}
The artifact contains 81 theorem and lemma declarations, all proved: 38 in \texttt{Core.lean}, 12 in \texttt{Ogd.lean}, 9 in \texttt{Hedge.lean}, 6 in \texttt{BoxD.lean}, and 16 in \texttt{Coupled.lean}.
The formalized chain covers the pull identity, the eigenbasis decoupling, the risk-interpolation and calibration identities, the minimax core, the lag-invariance step, the finite-cadence refresh chain (rate identity, fixed point, geometric closed form, contraction, variance identities, and the lower endpoint $\Psi \ge (1-s)^2$), the moving-frame reduction and cycle-average assembly, the subgradient and box-projection steps, the Zinkevich and Hedge regret bounds, the window optimum and second-derivative positivity, the spectral invariance under orthogonal conjugation and the conditioning bound, the deterministic-equivalent collapse, and the coupled contraction chain, together with the exact effective-dimension decomposition of the window coefficient.

The probabilistic concentration step, the deterministic-equivalent convergence of Theorem~\ref{thm:rmt}, and the Davis--Kahan bound are pen-and-paper.
The Lean declarations take the resulting recursions as hypotheses; Table~\ref{tab:lean-map} marks each row as formalized or pen-and-paper.

No declaration introduces a custom axiom: the sources contain no \texttt{axiom}, \texttt{sorry}, or \texttt{unsafe} construct, and the headline theorems' axiom dependencies are logged in \texttt{theory/lean/axioms.log}, regenerated by \texttt{theory/lean/print\_axioms.sh}, all resting only on \texttt{propext}, \texttt{Classical.choice}, and \texttt{Quot.sound}.

\subsection{Build Recipe}
{\sloppy The project ships its own \texttt{theory/lean/lakefile.toml}, \texttt{lean-toolchain}, \texttt{lake-manifest.json}, and root import module; \texttt{theory/lean/build\_ogd.sh} copies the sources and configuration into the lake project (default \texttt{\$HOME/SparTheorems}) and builds:\par}
\begin{quote}
\ttfamily
bash theory/lean/build\_ogd.sh\\
bash theory/lean/print\_axioms.sh
\end{quote}
On first use the script resolves Mathlib with \texttt{lake update} and fetches the pre-built cache with \texttt{lake exe cache get}, so no Mathlib compilation is required.

